% ============================================================
%  R. Nielsen, Forelæsninger over „Philosophisk Propædeutik“
%  fra Universitetsaaret 1860--61 (Kjøbenhavn 1862, 482 pp.)
%  "Lectures on 'Philosophical Propaedeutics' from the
%   Academic Year 1860--61"
%  TRANSLATION (English) — from the verified Danish transcription.tex,
%  which is COMPLETE (pp. 1--482, check.py clean).
%
%  Method: prose translation of the whole book, built batchwise.
%  The `% --- p. N ---` page-break comments are carried over from
%  transcription.tex verbatim so the two files stay aligned line for line;
%  chapter/section/subsection heads and \addcontentsline entries mirror the
%  transcription 1:1. Unfilled spans carry
%      % [text to be added: pp. X--Y]
%  markers in reading order. See ../../../TRANSLATION-PLAYBOOK.md.
%
%  HOUSE STYLE / KEY TERMS. Settled against the sibling Nielsen and Brøchner
%  translations in this repository, so that a reader moving between them
%  meets the same English for the same Danish:
%    den sunde Forstand      -> "sound understanding"      (in quotes, as printed)
%    den sunde Menneskeforstand -> "sound human understanding"
%    Anskuelse               -> intuition       (Anschauung)
%    Forestilling            -> representation  (Vorstellung)
%    sandselig / Sandsning   -> sensuous / sensation
%    Erkjendelse             -> knowledge, cognition
%    Tilværelse / Existens   -> existence
%    Væren / det Værende     -> being / that which is
%    Væsen                   -> essence
%    Ophæve / Ophævelse      -> sublate / sublation (Aufheben)
%    Modsigelseslære         -> Doctrine of Contradiction
%    Mulighedslære           -> Doctrine of Possibility
%    Tvetydighedslære        -> doctrine of ambiguity
%    Grundsætning            -> fundamental principle / principle
%    Sætning                 -> proposition (occasionally "principle")
%    Skarpsindighed / Spidsfindighed -> acumen / over-subtlety
%    Bestemmelse             -> determination
%    Momenter                -> moments
%    dialektisere            -> dialecticize (Nielsen's own coinage)
%    Idee                    -> Idea (capitalized where the Hegelian item is meant)
%  Hegelian capitals: Idea, Concept, Being, Spirit, Nature take a capital where
%  Nielsen means the technical item; otherwise nouns are lower-cased, despite
%  Danish orthography capitalizing every noun.
%
%  EMPHASIS. The transcription distinguishes THREE printed devices —
%  \emph{} (letterspacing/Sperrsatz), \textit{} (antiqua against Fraktur,
%  i.e. Latin/foreign words and italicized names), \textbf{} (bold Fraktur,
%  used for run-in heads and for the opening question of a section).
%  All three are carried over unchanged, span for span.
%
%  QUOTES. Danish „…“ -> ``…''; a quote inside a quote -> `…'.
%  ONE DELIBERATE DEPARTURE: Nielsen frequently runs a sample proposition into
%  his sentence unquoted, relying on the Fraktur/roman contrast and on Danish
%  word order to set it off ("saasom: Sort er ikke hvidt, Tykt er ikke tyndt";
%  "Sætte vi ubetinget: Dyr er Dyr"). English has neither resource, and the
%  sentences become unreadable without a mark. Such propositions are therefore
%  put in quotation marks here even where the print has none. This is the only
%  systematic source of the English quote count exceeding the Danish; the
%  running balance is still expected to be even.
%  Greek is copied verbatim from the Danish (textalpha loaded); Latin tags
%  stay Latin in \textit{}; the German and French quotations of the later
%  parts are to be given verbatim followed by a bracketed English rendering.
%
%  FOOTNOTES. Nielsen's own notes keep the printed *) / **) marks, driven by
%  the footnote counter with \setcounter{footnote}{0} after the page marker of
%  every printed page carrying a note, exactly as in transcription.tex.
%  TRANSLATOR'S notes are a separate apparatus: \tnote{...} sets a dagger and
%  does not disturb the author's counter. They are used sparingly, at the FIRST
%  occurrence of a technical term whose 19th-century sense diverges from
%  present-day English, and for the book's internal cross-references.
%
%  CROSS-REFERENCES. Nielsen cites his own earlier books by short title:
%    "Propæd." / "Propædeut."   = Philosophisk Propædeutik i Grundtræk (1857)
%    "Philos. og Mathem."       = Philosophie og Mathematik (1857 lecture)
%    "Mathem. og Dialekt."      = Mathematik og Dialektik
%  These are left as he gives them and glossed once, at the first occurrence.
%
%  PRINTER'S DEFECTS in the Danish are logged at their site in the
%  transcription and cannot be reproduced in English; where one affects the
%  sense, a % comment marks the spot here too.
%
%  STATUS: IN PROGRESS.
%    DONE — Forord; I. The Concept of Philosophy, pp. 1--29 and the dialogue's
%           tail on p. 30 (A. Doctrine of Contradiction; B. Doctrine of
%           Possibility; C. Dialectic, with the two dialogues "Something and
%           Other" and "Concept and Existence").
%           II. Philosophy and the Particular Sciences: General Remarks,
%           pp. 30--38 (philosophy's relation to the subject-sciences);
%           II.A, Philosophy's Relation to Mathematics — the section opening
%           p. 39 and the whole of a) The Concept of Infinity, pp. 40--45.
%    NEXT — % [text to be added: pp. 45--49] — II.A b) The Problem of Infinity.
%
%  MATHEMATICS. From p. 39 on the book sets equations inline in $...$. These are
%  carried over CHARACTER FOR CHARACTER from transcription.tex, including the
%  print's own inconsistencies (e.g. the multiplication dot after the first
%  \infty on p. 44, present in some equations and absent in others). Never
%  normalise the notation. The batch check is: extract every $...$ chunk from
%  both files, normalise whitespace, and require the two lists to be identical.
% ============================================================
\documentclass[12pt,a4paper]{book}
\usepackage[utf8]{inputenc}
\usepackage[T1]{fontenc}
\usepackage{libertinus}
\usepackage{libertinust1math}
\usepackage{textalpha}   % Greek typed directly (α) β) γ) subdivision marks)
\usepackage{graphicx}    % \reflectbox, for the old »ɔ:« id-est mark
\usepackage{tikz}        % the affinity/branching schemes, pp. 95, 96, 389
\usetikzlibrary{decorations.pathreplacing}

% Double-decomposition scheme, as printed on pp. 95 and 96 — kept identical to
% the transcription's macro so the later batches can reuse it unchanged.
\newcommand{\decompscheme}[6]{%
  \begin{center}
  \begin{tikzpicture}[baseline,every node/.style={inner sep=1.5pt,outer sep=2pt}]
    \node[anchor=west] (r1) at (0,2.55) {#1};
    \node[anchor=west] (r2) at (0,1.70) {#2};
    \node[anchor=west] (r3) at (0,0.85) {#3};
    \node[anchor=west] (r4) at (0,0.00) {#4};
    \node[anchor=west] (p1) at (5.4,2.125) {#5};
    \node[anchor=west] (p2) at (5.4,0.425) {#6};
    \draw[decorate,decoration={brace,amplitude=5pt}]
         (-0.22,2.85) -- (-0.22,1.40);
    \draw[decorate,decoration={brace,amplitude=5pt}]
         (-0.22,1.15) -- (-0.22,-0.30);
    \draw (r1.east) -- (p1.west);
    \draw (r3.east) -- (p1.west);
    \draw (r2.east) -- (p2.west);
    \draw (r4.east) -- (p2.west);
  \end{tikzpicture}
  \end{center}}
\DeclareUnicodeCharacter{0254}{\reflectbox{c}}  % ɔ (U+0254) = reversed c
\DeclareUnicodeCharacter{0186}{\reflectbox{C}}  % Ɔ, for CIƆ IƆCLXXVII on p. 425

\usepackage[english]{babel}
\usepackage[protrusion=true,expansion=true]{microtype}
\usepackage{setspace}
\usepackage[margin=1.2in]{geometry}
\usepackage{enumitem}
\usepackage{fancyhdr}
\setlength{\headheight}{28pt}
\addtolength{\topmargin}{-13.5pt}

% Nielsen's own footnotes, as printed: *) for a page's first note, **) for its
% second. The counter is reset by an explicit \setcounter{footnote}{0} placed
% after the % --- p. N --- marker of EVERY printed page that carries a note,
% exactly as in transcription.tex.
\renewcommand{\thefootnote}{\ifcase\value{footnote}\or *)\or **)\or ***)\fi}

% Translator's notes: a dagger, and the author's counter left untouched.
\newcommand{\tnote}[1]{%
  \begingroup
  \renewcommand{\thefootnote}{\textdagger}%
  \footnote{\emph{Translator's note.} #1}%
  \addtocounter{footnote}{-1}%
  \endgroup}

\usepackage{hyperref}

\hypersetup{
  pdftitle={Lectures on ``Philosophical Propaedeutics'' from the Academic Year 1860--61},
  pdfauthor={R. Nielsen},
  hidelinks
}

\pagestyle{fancy}
\fancyhf{}
\fancyhead[LE,RO]{\thepage}
\fancyhead[RE]{\textit{Lectures on ``Philosophical Propaedeutics''}}
\fancyhead[LO]{\textit{\leftmark}}

\setstretch{1.3}

\title{\textbf{Lectures}\\[4pt]
  \large on\\[4pt]
  \textbf{\large ``Philosophical Propaedeutics''}\\[6pt]
  \large from the Academic Year 1860--61\\[10pt]
  \normalsize translated from the Danish by H.\ Halvorson}
\author{R.\ Nielsen}
\date{Copenhagen: Published by the Gyldendal Bookshop (F.\ Hegel), 1862}

\begin{document}
\maketitle
\thispagestyle{empty}
\newpage

\tableofcontents
\newpage

%% ============================================================
%%  TRANSLATOR'S NOTE ON TERMINOLOGY  (not in the Danish)
%% ============================================================
\chapter*{Translator's Note on Terminology}
\addcontentsline{toc}{chapter}{Translator's Note on Terminology}
\label{ch:tnote}

Nielsen lectured in 1860--61 in a philosophical Danish that had been formed on
German models, and much of his vocabulary is a set of Danish calques on Kantian
and Hegelian terms. A few of these words have modern English cognates that mean
something quite different, and a reader who takes them at face value will be
misled. The renderings below are used consistently throughout; a dagger
footnote marks the first occurrence of each in the text.

\begin{description}[leftmargin=2.2em,style=nextline,itemsep=2pt]

\item[\emph{den sunde Forstand} --- ``sound understanding''] Literally ``the
healthy understanding''; the standing Danish equivalent of German
\textit{der gesunde Menschenverstand}, and close to English ``common sense.''
But \emph{Forstand} is \textit{Verstand}, the understanding as opposed to
reason, and Nielsen's whole argument turns on its being a \emph{faculty} that
can be shown to contradict itself --- not a stock of shared opinion. He almost
always sets it in quotation marks, and those are preserved. The fuller form
\emph{den sunde Menneskeforstand} is given as ``sound human understanding.''

\item[\emph{Anskuelse} --- ``intuition''] \textit{Anschauung}: the immediate
presentation of an object to sense, in Kant's technical sense. Never
``opinion'' or ``view,'' which is what the cognate suggests in modern Danish.

\item[\emph{Forestilling} --- ``representation''] \textit{Vorstellung}: the
mental image or presentation, standing between sensation and concept. Nielsen's
argument at pp.~10--11 depends on the contrast between what can be
\emph{thought} and what can be \emph{represented} --- a round square can be
thought, he says, but not represented.

\item[\emph{sandselig, Sandsning, Sandseindtryk} --- ``sensuous, sensation,
sense-impression''] \textit{sinnlich}: pertaining to the senses. ``Sensuous''
here carries none of the modern connotation of ``sensual.''

\item[\emph{Skarpsindighed} / \emph{Spidsfindighed} --- ``acumen'' /
``over-subtlety''] A matched pair in Danish (\emph{skarp} = sharp,
\emph{spids} = pointed) that English cannot match morphologically. The first is
the virtue of exact analysis; the second is its vice --- subtlety that dazzles
by distinctions only half true. Distinguishing them is the business of
Part~I.A, and Nielsen insists the same utterance may be either.

\item[\emph{Menen} / \emph{Viden} --- ``opining'' / ``knowing''] The Platonic
$\delta\acute{o}\xi\alpha$/$\dot{\epsilon}\pi\iota\sigma\tau\acute{\eta}\mu\eta$
distinction. Nielsen puns on it at p.~22 (``the question is not of opining but
of knowing --- still, what is your opinion?'').

\item[\emph{Ophæve, Ophævelse} --- ``sublate, sublation''] \textit{aufheben}:
at once to cancel, to preserve and to raise up. Where the context is plainly
non-technical the ordinary ``cancel'' or ``annul'' is used instead.

\item[\emph{Tvetydig} / \emph{Tvesidig} --- ``ambiguous'' / ``two-sided'']
Another pair the English breaks: literally ``two-meaninged'' and
``two-sided.'' Nielsen's definition of dialectic at p.~15 is that it is the
science which discovers the two-\emph{sided} in the two-\emph{meaninged}.

\item[\emph{Grund} / \emph{Følge} --- ``ground'' / ``consequent'']
\textit{Grund}/\textit{Folge}, the logical pair, not cause and effect.

\item[\emph{dialektisere} --- ``dialecticize''] Nielsen's own verb, and
deliberately ungainly in Danish too: to derive one thing from another by
dialectic. It is kept rather than paraphrased, since the dialogue of pp.~24--28
turns on the question of what can and cannot be so derived.

\item[\emph{Grundsætning} / \emph{Sætning} --- ``fundamental principle'' /
``proposition''] \emph{Sætning} does duty for both ``sentence,''
``proposition'' and ``theorem''; where Nielsen means one of the laws of
thought, ``principle'' is used.

\item[\emph{diversi respectus}] Left in Latin, as Nielsen leaves it: the
scholastic maxim \textit{diversi respectus tollunt omnem contradictionem} ---
different respects remove every contradiction. It is the hinge of the first
dialogue, and translating it would blunt the joke that B.\ has never troubled
to study what a \emph{respectus} is.

\end{description}

\noindent
Nielsen cites three of his own earlier books by short title, without further
identification: \textit{Propæd.} or \textit{Propædeut.} is
\textit{Philosophisk Propædeutik i Grundtræk} (Copenhagen 1857);
\textit{Philos.\ og Mathem.} is \textit{Philosophie og Mathematik};
\textit{Mathem.\ og Dialekt.} is \textit{Mathematik og Dialektik}. These are
left in Danish, as he gives them. All page references in such citations are to
the Danish originals.

\bigskip
\noindent
Nielsen's own footnotes are marked, as in the print, \hbox{*)} and \hbox{**)};
translator's notes are marked with a dagger and open with the words
\emph{Translator's note}.

\newpage
%% ============================================================
%%  PREFACE  (Forord; unpaginated, PDF 9--10)
%% ============================================================
\chapter*{Preface}
\addcontentsline{toc}{chapter}{Preface}
\label{ch:forord}

% --- Forord, PDF 9 ---
Should the present work attract any attention, the interest taken in it will
doubtless be of two kinds. With the question of the possible worth of its
contents, most readers, keeping to the letter of the title, will surely join a
second and more practical question: \emph{on what principles, and to what
extent, was the philosophical propaedeutic here delivered to the students?}

As regards this latter question, it will not be enough to remark merely that
the motives, leading ideas and whole design of the present exposition
correspond to the oral lectures, but that in matters of detail a good deal has
been altered, something added, something else omitted, and so forth; the reader
will think himself entitled to a clearer account.

It is therefore expressly noted that, apart from a few less considerable
editorial changes, the whole of the first part, pp.~1--30, is a
straightforward reproduction of the lectures as delivered. The same holds of
the introductory piece of the second part, \emph{General Remarks}, pp.~30--39;
of Philosophy's Relation to Mathematics, pp.~39--63; of Philosophy's Relation
to Chemistry, pp.~81--100 and pp.~135--177;
% --- Forord, PDF 10 ---
and of Philosophy's Relation to Physiology, pp.~177--180, pp.~199--288,
pp.~386--437, pp.~475--482.

In the presentation of the remaining portions, by contrast, considerable
changes have been made. The section on the psychophysical investigations was
not touched upon at all in the oral lectures; sharper and further-reaching
analyses have been introduced here and there; and it goes without saying that
by far the greater part of the contents of the notes was added in the course of
the revision.

Since the effort to give the treatment of the physiological problems in
particular that fullness which investigations of so intricate a nature can
hardly do without brought with it a considerable expansion of the second part,
the third part --- Philosophy and Academic Study --- has been confined to a
bare statement of its moments, and is appended as a separate part only in order
to exhibit the course of the whole.

\noindent 11 June 1862.

\begin{flushright}
\textbf{R.~Nielsen.}
\end{flushright}

% (double rule below the signature, closing the Forord as printed)

%% ============================================================
%%  I. THE CONCEPT OF PHILOSOPHY  (pp. 1--29)
%% ============================================================
\chapter*{I.\\ The Concept of Philosophy.}
\addcontentsline{toc}{chapter}{I. The Concept of Philosophy}

% --- p. 1 ---
% (chapter opening; no header numeral on this page, signature "1" at foot)
\section*{A.\\ The Doctrine of Contradiction.}
\addcontentsline{toc}{section}{A. The Doctrine of Contradiction}

\textbf{What is philosophy?} ``A science full of contradictions!'' However much
this definition may seem, viewed from one side, designed to pass sentence on
philosophy altogether, it is nonetheless well founded, both in the history of
the science and in the nature of the case. For any other science such an
admission would be annihilating; for philosophy it furnishes an essential
starting point, a dependable fulcrum, a firm and unshakable foundation. That
philosophy must contain contradictions is a matter of course, since it is
precisely the science which everywhere discovers and exhibits contradictions
--- in existence itself no less than in sound understanding's\tnote{%
\emph{Den sunde Forstand}: the Danish equivalent of German \textit{der gesunde
Menschenverstand}, and close to English ``common sense.'' But \emph{Forstand}
is \textit{Verstand}, the understanding as against reason, and Nielsen's
argument turns on its being a faculty that can be shown to contradict itself.
He sets it in quotation marks almost everywhere, and those are kept. See the
Translator's Note on Terminology.} grasp of existence (Propæd.\ pp.~40--44)%
\tnote{Nielsen's short titles for his own earlier books, cited throughout
without further identification: \textit{Propæd.} (also \textit{Propædeut.}) is
his \textit{Philosophisk Propædeutik i Grundtræk} (Copenhagen 1857);
\textit{Philos.\ og Mathem.} is \textit{Philosophie og Mathematik};
\textit{Mathem.\ og Dialekt.} is \textit{Mathematik og Dialektik}. Page numbers
refer to the Danish originals.} --- and it is just in this respect that
philosophy is to be understood as a \emph{general doctrine of contradiction}.

As a doctrine of contradiction, however, philosophy must not merely discover
and exhibit contradictions; it must also test them, for there are differences
among contradictions. Some are superficial, unnecessary, easily avoided; others
are insistent and not to be waved aside. Some creep in through thoughtlessness,
others
% --- p. 2 ---
are manufactured by over-subtlety, others again brought to light by
acumen.\tnote{\emph{Spidsfindighed} and \emph{Skarpsindighed} form a matched
pair in Danish (\emph{spids} = pointed, \emph{skarp} = sharp) that English
cannot match. The first names the vice of subtlety that dazzles by
half-true distinctions; the second the virtue of exact analysis. Distinguishing
them is the business of this section, and Nielsen insists that one and the same
utterance may be either.} Some contradictions play about as in a game of
thought; others mark essential knots, difficult problems not easily untied.

The following examples may serve to show how necessary it is, where
contradictions are concerned, to distinguish the superficial from the
essential, the acute from the over-subtle.

\textbf{The Liar} (ψευδόμενος). This contradiction is put as follows: when
someone says that he is lying, is he lying in that very statement, or is he
telling the truth? The answer demanded is either yes or no. ``Yes, he is
lying!'' --- And he says so himself, calls the thing by its right name, says
what is true, and hence is not lying. ``No, he is not lying, he is speaking the
truth!'' --- Then one must believe what he says; but what he says is precisely
that he is lying; and if it is true that he is lying, then he is not speaking
the truth.

That there is a contradiction here, sound understanding can readily grasp; but
what sort of contradiction is it, and where does the fault lie?

\textbf{Equal Magnitudes.} That a magnitude is equal to itself, and that two,
three~.\,.\ magnitudes are equal to one another --- in this sound
understanding sees no contradiction. But if it were asserted that two
magnitudes cannot possibly be equal to one and the same third, and hence equal
to each other, because that would be a contradiction, sound understanding would
feel itself affronted. Is there then really any contradiction whatever in the
proposition just cited? To see what is actually being spoken of here, one must
expressly note that the talk is not of things that have magnitude but of
magnitudes. If the talk were not of magnitudes themselves but of things that
have magnitude --- an apple, say, a triangle, an hour, and so on --- there
would certainly not be a trace of contradiction. That two triangles, each of
which is equal to one and the same
% --- p. 3 ---
third, are equal to one another, goes without saying. So long as we have to do
with concrete magnitudes, difference and plurality can coexist perfectly well
with equality; but when it is pure magnitudes that are in question, the matter
stands otherwise. A plurality of magnitudes necessarily presupposes that the
one can be distinguished from the other: without mutual difference, no
plurality. But by what should pure magnitudes be distinguished from one another
if not by size? If the magnitudes are equal to one another, then every glimmer
of difference has vanished, and with it the plurality: to say that several
magnitudes are equal to one another is therefore to say, in other words, that
they are not several but one and the same magnitude, which, however many times
it be repeated, is not thereby multiplied but remains the very same.

Should ``sound understanding'' wish to reject this consequence because there is
something offensive about it, that can hardly be done out of hand by a mere
dictum of authority; a sensible rejection must set about showing wherein the
difficulty consists and how it is to be removed. If one and the same magnitude
can, by being named twenty times, become twenty magnitudes equal to one
another, then one and the same species must likewise be able, by being named
twenty times, to become twenty species; but if it is not enough to repeat the
same magnitude --- if in repeating it a difference must also be thought ---
then where is that difference to come from?

\textbf{The Instant.}\tnote{\emph{Øieblikket}: literally ``the eye-blink,'' the
smallest unit of time. ``Instant'' is used rather than ``moment'' throughout,
to keep the word clear of the Hegelian \emph{Moment} that Nielsen also
employs.} If the old question --- when did Socrates die? --- is answered by
pointing to the year of his death, sound understanding will hardly suspect any
contradiction in that. Unless indeed it be that Socrates has no year of death
at all, inasmuch as he did not take a year over dying: he was hale and healthy
through the whole of that last year, until he drank from the cup. But if, in
accordance with a well-known and warranted usage, one means by his year of
death the year in which the day of his
% --- p. 4 ---
death fell, then there is no contradiction here. If it be objected that he did
not take a whole day over dying either --- well then, let one think of the hour
of his death; and if one will go to the very limit, of the instant of his
death. So then: when did Socrates die? ``At the instant he ceased to live!''
But this answer, which tells us only that he died when he died, is no answer at
all. The question concerns the point of time in Socrates' existence at which
his death occurred. To specify this point of time, four cases may be conceived:
either he died α) while he was still living, or β) when he was already dead, or
γ) when he was partly living and partly dead, or δ) when he was neither living
nor dead. But α) to die while one is still living, to be a living dead man, is
surely a contradiction. And to die when β) one is already dead is surely a
contradiction too. The instant at which Socrates died must therefore have been
the instant when, lying on the border between life and death, he was in such a
condition that one part of him --- his feet, say --- was already dead while
another part was still alive. By this explanation, however, the question is not
answered but only pushed further back; for the question then is: when did the
already dead part die? If the answer is that the already dead part died when it
was still partly alive and partly dead, then the contradiction concealed in
case γ) becomes obvious. And to conceive Socrates in a condition in which he
was neither living nor dead is, finally, likewise to conceive a contradiction
(δ); for how long does a living thing live? Until it dies!

\emph{``A philosophy that stands in open contradiction with `sound
understanding' is false; a philosophy that conforms itself immediately to sound
understanding is superficial. Philosophy does not wish arbitrarily to
contradict the natural understanding, but only to show that so-called sound
sense, without having
% --- p. 5 ---
any inkling of it, incessantly contradicts itself.''} (Propæd.\ p.~41).

In so far as contradictions of these various kinds develop in the course of
thinking, and their origin is viewed from the subjective side, the
psychological share that reflection has in the contradictions offensive to
``sound sense'' must be determined more closely. The aberrations belonging here
may be briefly characterized as the \emph{eristic}, the \emph{over-subtle} and
the \emph{sophistical}.

The striking mark of eristic is a one-sided bent for quarrel and strife, a
delight in shining everywhere by objections, in raising difficulties, in
tearing down what is generally accepted and stubbornly defending one's own
divergent claims. As for ancient philosophy, it is above all the Megarian
school, akin to the Eleatic, that has earned the name of eristic; but neither
can it be denied
% (a small stray mark stands between "kan" and "heller" in the Danish print;
%  see the transcription's note at this site)
that every distinctively developed form of philosophy has its polemical
element, and that eristic, for all its dark sides, is to that extent of
unmistakable importance in the development of the science. For insight and
knowledge it is not enough that one form what is commonly called a reasonable
opinion about things, and, by combining sense-impressions with natural
judgment, be led to distinguish correct from incorrect, acceptable from
objectionable opinions; by means of opinions one attains only the probable, not
the true. Whoever would rise from mere opining to knowing\tnote{%
\emph{Menen} and \emph{Viden}, the Platonic pair
$\delta\acute{o}\xi\alpha$/$\dot{\epsilon}\pi\iota\sigma\tau\acute{\eta}\mu\eta$.
Nielsen returns to the pun at p.~22.} must arm himself against the deceptions
of sense, let thought test what has been accepted, and acknowledge as true only
what admits of being thought --- declaring untrue, without regard to the
multitude's otherwise unimpeachable opinions, everything that does not admit of
being thought. That philosophy, on these terms, must come into manifold
conflict with sound human understanding is evident. (Cf.\ Philos.\ og
Mathem.\ p.~12, note; Mathem.\ og Dialektik pp.~34--35, pp.~116--117,
pp.~131--32, etc.)
% --- p. 6 ---
Rigorous analysis of concepts demands acumen; but acumen that is one-sided and
unfruitful, dazzling by distinctions which are subtle yet only half true or
even untrue, yields over-subtlety. Easy as it is to state the difference
between acumen and over-subtlety so long as one keeps to generalities, it is
just as impossible to lay down a rule valid for every particular case; for one
and the same utterance may, seen from one side, look over-subtle, and yet, seen
from another, bear the stamp of sound and genuine acumen. The contradiction set
up in ``the Liar'' may serve as an example. If the task is to drive home the
truth that one should not undertake off-hand to answer every question, since
there are questions whose apparently clear and easily grasped sense proves on
closer inspection to be senseless, then ``the Liar'' may stand as a model of
genuine acumen. If, on the other hand, one guilelessly takes up the answer
demanded --- and does so in such a way that, watching with astonishment as the
affirmation turns into a denial and the denial back into an affirmation, one
ends by imagining oneself before some deep riddle, some exceedingly intricate
knot whose solution may perhaps shed a hitherto unknown light upon ``the
eternal truths'' --- then over-subtlety has outwitted sound understanding and
caught it in its web.

In secret league with over-subtlety and hair-splitting we find sophistry. The
sophist is not content to exhibit contradictions; he manufactures
contradictions, plays with contradictions, draws inferences, and dazzles by
false inferences. For it is one thing to bring contradictions into relief,
another to misuse them in the service of sophistic; how great the difference is
may be seen, for instance, from what was said above about ``the Instant.'' No
sensible person can regard it
% --- p. 7 ---
as sophistry to call attention to the contradictions upon which thought must
necessarily strike when it tries to give an account of the meaning of the
instant; only when it is a question of the use made of the
contradictions does sophistry become recognizable, by what is false, untenable
and dazzling in the inferences. Thus if anyone were to infer, from the
contradictions contained in the thought of the instant of death, the
unthinkability of death, and from the unthinkability of death its
impossibility, and were then to build upon that a new proof of the immortality
of the soul, fitted out with a peculiar acumen, we should have to consider him
a sophist.

% (section rule as printed)
\section*{B.\\ The Doctrine of Possibility.}
\addcontentsline{toc}{section}{B. The Doctrine of Possibility}

If philosophy as a doctrine of contradiction is not to dissolve into
over-subtlety and sophistry, it must always find itself able, from one side or
another, to master the contradictions. It must then be able to determine their
relation both to things and to our reflections upon things, both to objective
being and to subjective thinking; in seeing through the essence of
contradiction it must discover the marks of what is in itself thinkable and
what is in itself unthinkable, and thereby, so to speak, wring from
contradiction a dependable measure of the acceptable and the objectionable, the
reasonable and the absurd, the possible and the impossible: a thoroughgoing
doctrine of contradiction is inseparable from a \emph{general doctrine of
possibility}.

\textbf{What, then, is philosophy?} ``A doctrine of the possible.'' But since
philosophy cannot possibly be a doctrine of everything possible, this Wolffian
definition is satisfactory only on condition that the concept of possibility
--- its extension, its content, and its relation to contradiction --- be
exactly determined. The principle that everything is possible which does not
contain a contradiction (\textit{possibile est, quod non implicat
contradictionem}) --- which in turn presupposes that everything containing a
% --- p. 8 ---
contradiction is impossible --- will, on the strength of what has gone before,
hardly furnish a dependable guiding line, whether dogmatic or critical.

In the first place it is self-evident that the principle as stated specifies
possibility only formally, and is therefore unfit to orient us with respect to
the real possibilities. Even if for the time being we declare everything that
contains a contradiction impossible, and everything that does not contain a
contradiction possible, the concept of possibility so given will still be so
negative, so empty and unsettled, that it cannot be brought into any positive
connection whatever with the actual. To determine the possibilities that hold
in actuality, it is necessary to know actual things together with their
conditions and laws; and knowledge of actual things is reached only by the road
of experience. Without the correctives of experience, dogmatizing thought with
its empty concepts will be as readily tempted to set the real possibilities too
narrow bounds as to allow them too wide a play. Or who, back in the days when
Wolff was fixing the bounds and defining the possibilities, would have seen the
possibility that ships should be able to go both against wind and against
current; that carriages without horses should be able to run --- not like
horses, but swift as flying birds; or the possibility that a man in London
should be able to correspond with a man in Calcutta as quickly and easily as if
they were talking together? So long as philosophy, with its formal concepts and
definitions, attempts to exhaust existence and to set itself up as judge over
what is possible and impossible in actuality, so long will sound human
understanding have manifold occasion for the triumphant application of Hamlet's
words: \textit{There are more things in heaven and earth, Horatio, than are
dreamt of in your philosophy.}

But even if philosophy gives up trying to fathom the positive, empirical
content of the real possibilities, it may
% --- p. 9 ---
still contribute essentially to purifying and heightening our knowledge of the
possible, provided only that it can secure for us the negative side of
possibility by determining the impossible. Now if the impossible can be known
only as what is in itself unthinkable, and if the unthinkable must always be
recognizable by contradiction, then the fundamental question presses itself
upon us of its own accord: \emph{how, then, does the contradictory stand to the
thinkable?} Here we are at a decisive turning point.

In presupposing without more ado the unity of thinking and being, dogmatic
philosophy proceeds at the same stroke from the assumption that the
contradictory is unconditionally an unthinkable, and the unthinkable an
impossible. But is the contradictory really unthinkable in every respect?

\emph{``What does contradiction mean? An existing nothing is a contradiction; a
round square is a contradiction; a flying horse is a contradiction too. And yet
these contradictions are of different kinds.''} (Propæd.\ p.~43).

If what is contradictory in itself is straightway unthinkable, how then does
thought come to apprehend this unthinkable? How does it happen that we not only
understand what is meant by the word ``contradiction,'' but also see clearly
what lies in the essence of contradiction? If we set, say, ``something''
against ``nothing,'' ``that which is'' against ``that which is not,'' thought
finds it just as easy to grasp the negative as to grasp the positive; and if we
set the proposition ``that which is, is'' against the proposition ``that which
is, is not,'' one notices at once that a conflict between subject and predicate
is as evident to thought as an agreement. But if thought apprehends the
conflict as readily as the agreement, sees through an existing nothing as
clearly as an existing something, and so far finds the former as thinkable as
% --- p. 10 ---
the latter --- how then can what is in itself conflicting and contradictory
show itself to thought as an unthinkable, and this unthinkable coincide with
the impossible?

To extricate oneself from the difficulty, let one try letting representation
come to the support of thinking.\tnote{\emph{Forestilling} is
\textit{Vorstellung}, the mental presentation or image standing between
sensation and concept; \emph{Anskuelse}, just below, is
\textit{Anschauung}, intuition in Kant's technical sense --- never ``opinion''
or ``view.'' The whole argument of this page turns on the contrast between what
can be \emph{thought} and what can be \emph{represented}.} Pure, abstract
thinking, detached from all intuition and representation, does indeed grasp the
negative as readily as the positive, since the negative is precisely
abstraction's own abstract product; but when the light of intuition dawns, when
representation joins in, the opposition between the positive and the negative,
between something and nothing, becomes conspicuous. For something admits of
being intuited, whereas nothing, pure nothing, does not; to the positive there
answers a representation, but nothing --- the negative --- repels every
representation, and is for that reason in turn repelled and rejected by
representation.

How thinking, by being combined with representation, comes to apprehend the
contradictory as the unthinkable, and the unthinkable as the impossible, may be
seen from the example of the round square. The proposition ``the square is
round'' can indeed be thought, in so far as the predicate's conflict with the
subject can be thought; but the round square cannot be represented, since the
representation of the round displaces and effaces the representation of the
square. Thought, seeing that the round square makes a demand upon
representation which representation must reject, now brands the contradiction
as something objectionable, declares the round square an unthinkable, and this
unthinkable an impossible. That it is precisely intuition and representation
which in such cases must help thought to discover the unthinkable is seen not
only from propositions in whose components representation manifestly rules ---
``black is not white,'' ``thick is not thin'' --- but also from examples that do
not directly display the contradiction in their words and manner of expression,
as when one asks: is
% --- p. 11 ---
a regular hexahedron possible?\tnote{The example is puzzling as it stands,
since the regular hexahedron is simply the cube. Nielsen's point is presumably
that these are questions whose answer cannot be read off the words --- unlike
``black is not white'' --- so that intuition must be consulted before one can
say whether a contradiction is present at all.} is it possible for a line
parallel to an ellipse to be an ellipse? and so on.

If only we could rely on intuition and representation to furnish, always and in
every case, infallible marks of the thinkable and the unthinkable, we should at
least have something firm to hold to here; but now look at the example of the
flying horse! Natural sense surely feels that under the name ``flying horse'' a
contradiction lies hidden. But to whom is sound human understanding now to turn
to have the contradiction cleared up? To thinking! That is no help: pure
thinking thinks everything, the unthinkable included. To intuition! No help
either: intuition lives with the poet, and the poet rides Pegasus, and Pegasus
reconciles intuition with every winged horse. To a union of thinking and
intuition! But how can a union of two unreliables yield anything reliable? To
sense-impressions! Sense-impressions can indeed, in union with thinking, teach
us something about the possible, in so far as they teach us something about the
actual, of course; but what do they teach us about the merely possible, and
what could they possibly teach us about what is in itself contradictory and
impossible?

To determine the relation of possibility to contradiction, it is above all
necessary to grasp \emph{the principle of contradiction} --- its form, its
validity, its application.

In formal logic we find the propositions of identity and of contradiction set
out so separately that we are thereby supplied not merely with two fundamental
principles but, into the bargain, with two independent principles. If there is
really anything in this duality of principles, we should expect some account of
the difference of principle between them. But wherein is the difference
supposed to consist? To explain the proposition of identity --- \textit{A} is
what \textit{A} is --- one slips in underneath it the proposition of
contradiction: that \textit{A} is
% --- p. 12 ---
what \textit{A} is precisely because \textit{A} is not other, and cannot be
other, than what \textit{A} is. Now if the principle of contradiction has had
to do service in setting up pure identity --- and that, in the logical sense,
before it was even born --- then it is no more than fair that the principle of
identity, once born and the logical birth-pangs happily over, should in its
turn help the principle of contradiction to due recognition. Why, then, is no
thing what it is not? Answer: because every thing is what it is. And why is
every thing what it is? Answer: because no thing is what it is not. In other
words: when identity is asked after, one answers with contradiction; and when
contradiction is asked after, one answers with identity. On closer inspection
it must undeniably strike sound human understanding that there are here not
really two fundamental principles but only one.

Yet suppose that, in order so far as possible to avoid conflict with
traditional forms, we let formal logic keep its two fundamental principles
undiminished --- so that, for the sake of ``certainty'' and ``truth'' and
``determination,'' ``statuation,'' ``whatness,'' ``decision,'' ``unanimity''
and ``concordance,''\tnote{The list is ironic: Nielsen is parading the
scholastic-sounding coinages under which formal logic defends its two
principles. \emph{Hvadhed} is a Danish calque of \textit{quidditas},
``whatness''; \emph{Statuering} is the laying down of something as established;
\emph{Eenstemmighed} and \emph{Samstemmighed} are ``unanimity'' and
``concordance.''} --- we held identity literally apart for itself and
contradiction apart for itself: would such a holding-apart then perhaps help
sound understanding to a clear and distinct insight into the relation between
the contradictory and the impossible? At first glance it seems that some
difficulty still remains. For if every \textit{A}, provided only that it is
what it is, has by the principle of identity a claim to ``a valid is,'' then it
looks almost as though the senseless and the thoughtless could likewise make a
claim --- a valid claim --- to ``a valid is''; if every thing, provided only
that it is not what it is not, can satisfy the principle of unanimity and
concordance, why should the conflicting, the absurd, the self-contradictory and
the foolish
% --- p. 13 ---
not be able to satisfy the principle of unanimity and concordance as well?

But to overlook the principle of contradiction, still less to reject it, will
not do at all; on the contrary, when thought attempts to deny the validity of
the principle of contradiction, it is led by the very denial to confirm it. To
posit ``black is white'' instead of ``black is not white'' is not only possible
but reasonable and necessary if the corresponding contradiction is to be made
properly plain; yet in being branded in this way as a contradiction, the
contradiction vindicates the principle. The unthinkable, in this sense, is not
a something that thought simply cannot think; the unthinkable is the thought
which thinking itself condemns, in view of a purpose given within a given
context --- a thought that comes into being only to vanish, a thought in whose
sublation\tnote{\emph{Ophævelse}, Hegel's \textit{Aufhebung}: at once
cancelling, preserving and raising up. Where the context is plainly
non-technical, the ordinary ``cancel'' is used instead.} thinking recognizes
its own law and satisfies itself. Only when the unthinkable is thus
determinately set against the thinkable can the latter be called the logical
mark of the possible and the former that of the impossible. Contradiction,
then, logically regarded, cannot be abolished outright. Could the principle of
contradiction abolish contradiction once and for all, it would at the same
stroke render itself superfluous. But since thought cannot stir without
contradiction --- its sting and its spur --- stirring too, the principle stands
in force because the cooperation of contradiction stands in force; in order to
determine the thinkable, sift the possible and dismiss the impossible, the
principle of contradiction must incessantly contradict contradiction, must in
resolving contradiction uphold contradiction.

In the application of the principle of contradiction the difference between
critical and dogmatic philosophy comes into view. Agreeing with dogmatism in
setting contradiction outwardly hostile against thought and the movement of
thought, critical philosophy then diverges from the dogmatic in this: that in
the unthinkable it sees only the subjective side of the impossible, and not its
% --- p. 14 ---
objective side as well. When Kant attempts to think the origin of the world, he
strikes, quite rightly, upon contradiction; the same recurs when he tries to
think the essence of substantiality, of necessity and of freedom objectively.
The acute thinker, consistent from his own point of view, is here so far from
being rid of the contradictory that on the contrary, seeing contradiction tie
its knots and settle into antinomies that cannot be waved aside (Propæd.\
pp.~29--30; Mathem.\ og Dialekt.\ pp.~125--126), he declares it insuperable.
The objective essence of the world belongs therefore, according to Kant, to
what is for us unthinkable; but what does Kant infer from this? By no means
that the reality of an objective world should on that ground be held impossible
--- that is an inference which first arises with Fichte (Propæd.\ p.~33) ---
no, the author of ``Kritik der reinen Vernunft'' warns precisely against
precipitate inferences from the subjectively unthinkable to the objectively
impossible. The separation of thinking and being, the original characteristic
mark of critical philosophy, rests precisely on the presupposition that for the
human understanding there are essential, necessary, insoluble contradictions.

But are the Kantian antinomies really insoluble? That depends on whether one
chooses to grant contradiction a merely subjective significance, or a
significance at once subjective and objective. Such a choice cannot of course
be arbitrary; it must follow from the nature of the case. If, then, it turns
out that by denying the objective validity of contradiction, holding fast to
the critical separation of thinking and being, and so rejecting every attempt
to know ``\emph{the thing in itself},'' one by no means escapes contradiction
but rather entangles oneself in new and greater contradictions --- what then?

% (short centred rule after this paragraph, closing the discussion as
%  printed at the foot of p. 14; p. 15 opens the next section)

% I. The Concept of Philosophy continues; C. Dialectic region.
% --- p. 15 ---
\setcounter{footnote}{0}
\section*{C.\\ Dialectic.}
\addcontentsline{toc}{section}{C. Dialectic.}

% FOOTNOTE MARK: the first footnote of the printing (below) is marked *) both
% in the text and at the note, which is set off by a short rule; the note
% continues at the foot of p. 16 without a repeated mark.

What, then, is philosophy? ``A doctrine of ambiguity, a science of the
ambiguous.'' The correctness of this definition seems at least to find ample
confirmation in what has gone before. For although we did distinguish between
acumen and over-subtlety, we had to admit that one and the same thing could be
taken both as over-subtle and as not over-subtle; when it was a matter of
separating the thinkable from the unthinkable, what came out was that the
unthinkable was something that both could and yet could not be thought; and
when it finally came to a resolution of contradiction, the upshot was that
contradiction both could and could not be resolved. Should sound understanding
find such a predilection for the ambiguous suspicious, let it be noted that the
expression ``doctrine of ambiguity'' is itself ambiguous. While higher
instruction in making the clear obscure, the distinct veiled and the
intelligible unintelligible is always to be sought among the sophists,
philosophy must contrive to set the ambiguous in another light --- a light that
makes the ambiguity vanish. As the science which everywhere discovers the
two-sided in the ambiguous,\tnote{\emph{Tvetydig} and \emph{tvesidig}:
literally ``two-meaninged'' and ``two-sided.'' The definition depends on the
Danish pair, which English cannot reproduce.} and by means of contradiction
clarifies the doubling, philosophy is \emph{dialectic}\footnote{That the form
of philosophy is in its essence dialectical must not be misunderstood as though
the development of concepts had everywhere to display immediately the formal
antinomies by which logical transitions are motivated; on the contrary!
Philosophical exposition can and should also take up reporting, descriptive and
critical components. The only decisive point is that the elements communicated
in the form of report, description, critical discussion and the like should
everywhere show themselves, by the design of the whole and the prevailing
current of thought, to serve the dialectical. There is in the Hegelian method a
formal overstrain, a logical over-exertion, brought on chiefly by the endeavour
to have the exposition sparkle at every point with
% (the footnote text continues at the foot of p. 16:)
dialectical flashes of lightning. It is with dialectic's flashes of lightning,
and with the elements in which they are prepared and out of which they develop,
as it is with natural lightning: the lightning-fire must as a rule presuppose
clouds, but the clouds are not on that account lightning.}.

% --- p. 16 ---
The truth that philosophy by its higher nature is and must be dialectic was
categorically stated by Plato. (Cf.\ Propædeut.\ pp.~124--129; p.~150 and
elsewhere.) In Plato the presentation of the dialectical (διαλεκτική from
διαλέγεσθαι) has, as is well known, taken the form of conversation; and
connected with this is the fact that the dialectician is characterized
precisely as the one who in conversation knows how to ask and to answer rightly
(Τὸν δὲ ἐρωτᾷν καὶ ἀποκρίνεσθαι ἐπιστάμενον ἄλλό τι σὺ καλεῖς ἢ διαλεκτικόν;
% „ἄλλό“ with double accent so printed in the Danish
\textit{Krat.}\ 390 c.). That what essentially matters in dialectic is not the
conversation but the development of concepts conditioned by ``the art of
language'' goes without saying; but that the conversational form, when
everything irrelevant is kept out, is nonetheless a very natural form for the
beginning dialectician, may readily be shown by a few examples.

\textbf{Something and Other.}
\addcontentsline{toc}{subsection}{Something and Other}
A. If the categories ``something'' and ``other'' simply excluded one another,
there would be no dialectic; if they coincided without more ado, there would be
no dialectic either. The dialectical shows itself precisely in this, that the
two propositions ``something is what it is, not an other'' and ``something is
what it is not, hence an other'' not only cancel but also mutually presuppose
and ground one another. B. If the dialectical rests on so slippery a ground,
then all dialectic is false. A. How so? B. The claim you have set up stands in
open contradiction with sound human understanding, and \emph{a philosophy that
stands in open contradiction with `sound understanding' is false}. A. If what
you here call \emph{open contradiction} is to be dismissed out of hand by an
immediate verdict, philosophy can attach no weight to the ruling; for \emph{a
philosophy that conforms itself immediately
% --- p. 17 ---
to `sound understanding' is superficial}. But if ``sound understanding'' will
condescend to an inquiry, the inquiry can easily be got under way. B. It would
have to be the second proposition, under whose negative form some ambiguity or
other has hidden itself, and which might therefore become an object of inquiry;
in the first proposition, the proposition of identity, the pure proposition of
certainty (\textit{principium certitudinis}), there cannot be a trace of
ambiguity. A. But suppose it is after all the first proposition in which the
ambiguity, the indeterminacy, and so far the uncertainty, lies hidden --- what
then? B. Then it must be capable of proof. A. Is it altogether forbidden in
logic to say anything else of something than that it is something? B. By no
means; for how, by merely repeating ``something is something,'' should one ever
determine \emph{what} something is? Logic, however, which has no fondness for
ambiguities, draws a definite distinction, as easy to grasp as it is
unmistakable, between determining something by means of an other and saying of
something something other than what it is. A. So it is permitted to determine
what something is by determining what it is not? B. Not only permitted but
necessary. A. Then it is necessary that in the determination of something a
determination of an other be contained. B. That too may be granted --- but on
the express condition that something, precisely by that determination, excludes
the other. A. Something, then, is~---? B. A thing different from an other.
A. And if now, on that condition, we exclude the other, something is~---?
B. A different thing.\tnote{\emph{Et Forskjelligt}, the substantivized neuter
adjective: ``a differing one,'' something whose whole character is to be
different. English has no natural form for it; ``a different thing'' is used
throughout the exchange, and the reader should hear the emphasis on
\emph{different} rather than on \emph{thing}.} A. But now the other, on its
side, is~---? B. A thing different from something. A. And when, likewise on
that condition, we let the other exclude something, we say that the other
is~---? B. A different thing. A. So then: something, when determined, is a
different thing; the other, when determined, is a different thing; something is
therefore what the other is, namely a different thing. B. This ambiguity is
easy enough to see through. Certainly
% --- p. 18 ---
something is a different thing, just as the other is; but different and
different are two things --- something is a different thing precisely by being
different from the other, and the other a different thing by being different
from something. Something and other are therefore still, as before, different
from one another. A. Certainly something and other are different from one
another; only a sophist would take it into his head to deny their difference.
What was to be proved here was only this: that their difference is grounded in
their identity. B. What is different is not identical, and what is identical is
not different: only by pure categories can one secure oneself both against
hidden ambiguities and against open contradictions.

A. But if I am not mistaken, you have already yourself, by your definite and
lucid answers, made valuable contributions toward sharpening our eye for the
dialectical unity of something and other. B. I should very much doubt it.
A. But you did stress that something can be determined only as a different
thing, and, for clarity's sake, as a thing different from an other? B. Yes,
that I did stress. A. Then you have at the same stroke expressly required that
something, in virtue of that difference, must be taken as an other than that
from which it is different. B. Of course. A. So on your view something is
other than the other, in opposition to the other --- and hence, after all, an
other. B. Your play on words leads nowhere; for when something is taken as an
other, it becomes another other --- namely the other that something itself is,
as against all the other that something is not. A. That I can understand; now
you are talking like a true dialectician. So you have now seen for yourself
that, with respect to the determination contained in something, there are two
kinds of other: an other that something itself is, as against another other
that something is not --- in other words, that something can exclude the other
only by itself being an other. B. Yes --- so long as one keeps to expressions
as general as ``something'' and ``other,''
% --- p. 19 ---
one can no doubt practise without much difficulty the game of thought of
turning something into other and other into something. But let us try it with
something determinate! Show me in definite examples that something really is
identical with other, and I shall grant you that there is ``sound
understanding'' in the dialectical. A. You take a leaf in a book to be
something determinate? B. Naturally. A. And the two sides of a leaf are surely
so opposed to one another that if we call the front a something, we must call
the back an other. B. Excellent! Now surely it will occur to no one to identify
the front with the back, and then maintain that the front is what the back is,
because both are sides, or because the front becomes the back when the leaf is
turned, and the back the front. A. But on this occasion too the identity is
surely as conspicuous as the difference? B. I can see no identity. A. What do
you think, then, about the leaf? B. That the leaf which has the two sides
consists precisely in the difference and opposition of the sides! For if the
two sides were turned into one side, we should get only a surface, not a leaf.
A. And if the two sides were separated in such a way that each came to belong
to a body of its own, we should get two surfaces but no leaf. B. Yet it makes
an essential difference whether I say ``in the leaf the front is united with
the back'' or ``the front is what the back is, the front is back.'' The first I
can grant, the second not. A. If the leaf that unites the two sides were a
third thing subsisting on its own over against the sides, we should certainly
not be entitled to take it as the existing unity of the sides; but since the
leaf, taken in separation from the sides, does not exist at all, the one side
must in the leaf show itself identical with the other. B. Are we then perhaps
to derive from this that the front itself is both front and back, and the back
in turn both back and front? A. That is a consequence you will hardly reject.
If we call the two sides α and β, it is
% --- p. 20 ---
evident that α, which in and for itself is neither front nor back, must by its
opposition to β be front in so far as β is back, and conversely; α is
therefore, in virtue of its identity with β, both front and back. B. That is a
peculiar extension of the concept of identity. What I have to object against
the concept of identity so extended I would rather hold back, however, until
you have produced one more example of a different sort.

A. What do you say to calling the greater a something and the lesser an other?
B. That example I find very suitable. A. And what is it, then, that we are to
prove? B. What else than that something is both a greater and a lesser, the
other both a lesser and a greater, and one and the same magnitude itself both
lesser and greater. And in truth nothing is easier than to prove all this. A
cursory comparison of three magnitudes, $\alpha > \beta > \gamma$, tells us at
once that α is a greater in so far as β is a lesser, and β a greater in so far
as γ is a lesser. But what follows from this? That β is in one respect a lesser
and in another respect a greater --- in that there is no contradiction; but if
anyone maintains that the magnitude β is at one and the same time and in one
and the same respect both greater and lesser than itself, he does fall into
contradiction, and not only with the natural understanding but with true logic.
That one and the same side of the same leaf can at one instant be front and at
another back, and at one and the same instant --- when, mark you, it is
regarded from two opposite points of view --- both front and back, contains no
contradiction, but a confirmation of the old rule \textit{diversi respectus
tollunt omnem contradictionem}. If, on the other hand, one ignores the
\textit{diversi respectus}, overlooking what both Plato and Aristotle, for all
the difference of their intellectual bent, so strictly inculcated --- that in
such inquiries what is asked about must each time be seen to be τὸ αὐτό, and
regarded κατὰ τὸ αὐτό,
% --- p. 21 ---
πρὸς τὸ αὐτό, ἐν τῷ αὐτῷ χρόνῳ --- then one may indeed dazzle the
inexperienced for a moment with piquant turns of identity and surprising
contradictions, but one will never convince the expert, never satisfy the
sensible thinking that inquires after truth. A. Weighty words, and important!
So now your ``sound understanding'' has almost got the better of
contradiction and finished off dialectic; I say \emph{almost}, for one small
point still remains --- a point we simply must have cleared up, and one you
will surely be able to clear up easily: it concerns \textit{diversi respectus},
and you have of course studied \textit{diversi respectus}. B. Studied? It has
never occurred to me that there was anything further to study in logic's
\textit{diversi respectus}; I have always supposed it enough to remark in a
general way that there are \textit{diversi respectus}. A. Then it is surely
necessary to distinguish between \textit{respectus}, since there are, after
all, \textit{diversi respectus}. Of the two sides of the same leaf --- to
return to one of our examples --- it undeniably holds in one respect that they
are the same, namely that both are sides, and in another respect that they
differ, namely that one is front and the other back; in one respect the same,
that both are, say, written on, in another respect different, that one is, say,
white and the other red, and so forth. Each of these different respects may in
turn give occasion for new ones; for although it holds of both sides in one
respect that both are written on, yet on this very point it also holds in
another respect that they differ, namely that the one is written in black ink
and the other in red. These doublings of respect can be continued indefinitely.
If thinking is not to lose itself in an unsurveyable tangle of all manner of
regards to all manner
% (see the transcription's note here: a mark above the n makes both OCR
%  witnesses read „allehaaude“ in the Danish; the reading was rejected)
of respects and remarks and considerations and reflections, logic must
necessarily distinguish not only between essential and inessential respects but
% --- p. 22 ---
also between respects that fall indifferently outside one another and respects
that are contained in one another and therefore follow from one another with
inner necessity. That one side of a leaf is written on does not of course
entail that the other must be so too; one can take account of the former
without taking account of the latter; the former can be without the latter's
therefore being, and no one but sophists would think of positing the latter as
identical with the former. It stands otherwise, however, with being front and
being back: here it is impossible to take account of the former without taking
account of the latter, for these two modes of existence are contained in one
and the same existent in such a way that the regard to the former actually
coincides with the regard to the latter --- the former is, despite the
difference, indeed in virtue of the difference, \emph{identical} with the
latter. True insight into the essence of identity, difference and contradiction
depends on analysing the \textit{diversi respectus} sufficiently.

B. But can any well-founded objection really be made against the proposition
\textit{A} is what \textit{A} is? A. Among others, the objection that the
proposition lacks a ground. B. But when the proposition is one with the
principle --- the principle of identity, the principle of truth, the principle
of certainty, \textit{principium certitudinis} --- and is of course
self-evident in every way, it needs no ground. A. Is it then immediately
evident from the proposition \emph{that} \textit{A} actually is? B. No, not
\emph{that} \textit{A} actually is, but \emph{what} \textit{A} is. A. Is it
immediately evident \emph{what} \textit{A} actually is? B. Not exactly
\emph{what} \textit{A} actually is, but that, \emph{if} \textit{A} is at all,
\emph{then} \textit{A} is \emph{what} \textit{A} is. A. But suppose \textit{A}
is not at all? B. Then it is still what it is. A. So in order to be what it is,
\textit{A} does not even need to be? B. You understand very well what I mean.
A. The question here is not of opining but of knowing;
% --- p. 23 ---
still --- what is your opinion? B. That even if \textit{A} is a non-being, it
is still what it is, namely a non-being. A. But non-being is surely the same as
what we otherwise call nothing, no thing? B. Naturally. A. So then, whether I
say ``no thing is what it is'' or ``every thing is what it is,'' it must at
once be evident to sound understanding --- by the principle of identity, of
course, the principle of truth, the principle of certainty, \textit{principium
certitudinis} --- that with both propositions I mean one and the same thing!

B. But can anything at all be adduced in defence of so self-contradictory a
proposition as ``something is what it is not''? A. Among other things this:
that a proposition which by its open contradiction so boldly challenges
thinking must awaken reflection and call for grounds, in that it demands a
refutation. B. But how should one refute by grounds a proposition that
immediately refutes itself? A. The proposition ``something is what it is not''
by no means refutes itself immediately. B. If what is unthinkable, absurd and
senseless in this proposition is not self-evident, then what is self-evident,
and what is sound human understanding? A. Sound human understanding by no means
sees anything absolutely absurd in something's being able to be what it is not.
B. So sound human understanding is supposed to be able to conceive that
\emph{something} --- mark you, at a given time and in a quite definite sense
--- \emph{is what it} --- at the same time and in the same sense --- \emph{is
not}; no, that really is too much. A. Yes, that really is too much; but so
absurd a demand is made by the proposition \emph{something is} --- mark you,
when its existence is grounded --- \emph{what it is not} only when the
proposition is held one-sidedly in superficial, abstract immediacy. B. Well,
that is another matter, if your defence rests on an interpretation.
% --- p. 24 ---
A. And another matter too, presumably, if your refutation rests on an
interpretation. B. Then prove to us, by means of your dialectical
interpretations, that one and the same something can at one and the same time
and in one and the same respect be both what it is and what it is not, and that
one and the same magnitude can at once be both greater and lesser than itself.
A. Such proofs belong to sophistic, not to dialectic. B. But to prove that
something can in one respect be what in another respect it is not, or that one
and the same magnitude can in one relation be a lesser and in another a
greater, nothing is needed but \textit{diversi respectus}, no dialectic at all.
A. But to prove that something, though different from the other, is
nevertheless essentially one with its other, because the being of both has the
same source, the same root; to prove that the same magnitude is on one and the
same ground both a greater and a lesser, since the two relations, the two
respects, have one and the same point at which they meet, namely being a
magnitude; to prove that the identity
% (the Danish prints „Identiten“ for „Identiteten“ here; see the transcription)
of something and other, of the greater and the lesser, contains precisely the
ground of something's being in existence different from the other, and the
greater different from the lesser --- for that, dialectic is needed; for this
dialectic can prove, but dialectic alone.

\textbf{Concept and Existence.}
\addcontentsline{toc}{subsection}{Concept and Existence}
C. A dialectical development of the inner connection of the \textit{a priori}
concepts, and of their transition into one another, is at the present standpoint
of the science no doubt to be acknowledged on the whole; but it is another
question whether, because the dialectical plays a considerable part in
philosophy, one is on that ground entitled not merely to reproach formal logic
with not being dialectical, but even --- on the strength of a Platonic
expression which can always be taken in more than one sense --- to transform
the whole of philosophy into sheer dialectic. A. Had you followed the preceding
inquiry attentively, you would,
% --- p. 25 ---
as to the first point, easily have convinced yourself that there was no thought
here of delivering a detailed critique of logical formalism, but only of
rebuffing the self-sufficiency with which one not seldom allows oneself,
invoking ``sound human understanding,'' to consign every logical objection to
the immediate validity of immediate first principles to the realm of fancy and
sophistry; as to the second point, your misgivings could do with a clearer
statement. C. \textit{A priori} concepts --- such as something and other,
essence and semblance, ground and consequent, and the like --- can easily be
produced out of one another by the application of contradiction. Whoever posits
the ground already implies the consequent; the ground is therefore, at bottom,
both itself and the consequent.\tnote{Danish \emph{Grund} carries the whole
weight here: it is at once the logical ground (\textit{Grund} as against
\textit{Folge}, consequent), the ``bottom'' of a thing, and the ``reason why.''
Nielsen plays on all three, and the English can only follow at a distance.} But
since the ground becomes a ground proper only on account of the consequent, the
ground is consequent and the consequent ground. The same thing is not merely
ground in one respect and consequent in another; the two respects also coincide
in a single respect: to dialecticize\tnote{\emph{Dialektisere}, Nielsen's own
verb and deliberately ungainly in Danish too: to derive one thing out of
another by dialectic. It is kept rather than paraphrased, since the whole
dialogue turns on the question of what can and cannot be so derived.} the
concepts out of one another is, if I may use the expression, the same as to
dialecticize \textit{diversi respectus} into one another. As for the \textit{a
posteriori} concepts, the concepts of experience, it is impossible to
dialecticize existence out of the concept; or who would try to dialecticize the
actual horse out of the concept of a horse, or the actual snail out of the
concept of a snail? But if it is impossible to dialecticize the existent out of
its concept, then presumably it is for the same reason impossible to
dialecticize one concept of existence out of another --- the concept of a
horse, say, out of the concept of a fly, and this in turn out of the concept of
a snail. In the world of experience \textit{diversi respectus} fall outside one
another, and where \textit{diversi respectus} fall outside one another
dialectic plays a very subordinate part. A. But if you could not succeed in
dialecticizing existence out of the concept (cf.\ Propæd.\ p.~31), have you
never tried going the opposite way --- dialecticizing the concept out
% --- p. 26 ---
of its existence, as for instance dialecticizing your concepts of snails out of
their corresponding snail-existence? C. One does not dialecticize concepts of
experience out of existing objects; one \emph{abstracts} them. A. So there is
nothing one might call a dialectic of abstraction? C. Not so far as I know.
A. But the concepts abstracted from the objects must surely be, not only in our
consciousness, but also in the objects themselves? C. Naturally; for if the
concepts were not in the objects, they could not be abstracted from the
objects. A. But the object from which the concept is abstracted is surely not
itself a concept? C. By no means. A. Yet the concepts, in so far as they are in
the objects, may be said to exist? C. That the concepts of experience exist in
the existing objects, no one can well doubt. A. So in one and the same
existence there are two existents --- the existing object, and the concept
existing in the object? C. No, it cannot be taken that way; it is not the
object itself but its material that is opposed to the concept. A. Then
presumably we get three existents: the existing material, the existing concept,
and that existence of the object in which both exist? C. No, one cannot take it
that way either. A. How, then, is it to be taken? C. The material by itself is
no separate existent, any more than the concept is; the material and the
concept are two sides --- two sides of existence --- of the same existing one,
the object. A. The existing one is, then, over against and in opposition to the
two sides of existence, the existent in and for itself? C. No, one cannot take
it that way either; for if we let abstraction remove the two sides of
existence, the existing unity will itself become an abstraction, a side of
existence that does not exist. A. But how, then, is it to be taken? C. So that
the existing unity is seen to exist in the two sides of existence, which
precisely thereby come to exist in inseparable unity. A. Then however many
parts,
% --- p. 27 ---
qualities and sides the existing object may otherwise have, we may now, on your
account, assume that these parts, with their sides, qualities and mutual
relations, are all equally material by virtue of the material and intellectual
by virtue of the concept; every side of existence thus becomes an existential
unity of two sides, the material and the intellectual. C. That we may quite
safely assume. A. But what is material in the object is surely something
sensuous,\tnote{\emph{Sandselig} = \textit{sinnlich}, pertaining to the senses;
it carries none of the modern connotation of ``sensual.''} which as sensuous
cannot be thought, and what is intellectual in the object something
non-sensuous, something thinkable that cannot be sensed? C. Perfectly correct.
A. Then in the existing unity the unthinkable is after all a thinkable, and the
non-sensuous a sensuous. C. That I do not understand! A. Then neither do you
understand the existing unity; for how should the thinkable in an existing
unity be able to penetrate the sensuous without the sensuous thereby becoming,
by virtue of the thinkable, itself thinkable, and the thinkable, by virtue of
the sensuous, sensuous? C. It almost looks as though there were something in
what you say. A. The concept of existence, then, is a system not only of
intellectual but also of material determinations --- a system of intellectual
sense-determinations, of sensuous intellectual determinations. C. Yes, so it
is. A. But then the concept of existence betrays a most striking contradiction.
C. What contradiction? A. That of demanding existence by setting itself against
existence. C. What does that mean? A. If the whole system of sensuous
intellectual determinations is posited in the intellectual form, then the
sensuous vanishes as a moment of reflection sublated in the thinkable --- that
is to say, the existent transforms itself into the concept of existence, but
existence goes out. If conversely the whole system of intellectual
sense-determinations is posited in the sensuous form, answering to immediate
sense-impressions, then the object steps into existence, but the concept of
existence goes out.\,.\,.\,. C. No, the concept of existence does not go out;
it is completed
% --- p. 28 ---
in existence: the existent is, by what we have said, nothing other than the
concept of existence complete in every respect. A. But was it not you who, by
what we have said, granted that the object from which the concept is abstracted
was not itself the concept? C. Yes, but I was thinking then of the abstract and
so far incomplete concept, not of the concept concrete and complete in itself.
A. So it is after all possible to dialecticize existence out of the concept,
once one has learned to dialecticize the concept out of existence? C. To prove
that in deed would here amount to showing it in definite examples.

A. Sound human understanding tells us that the horse is an animal, that the fly
is an animal, and the snail likewise; but how do such judgments actually stand
to ``the principle of identity and of contradiction''? C. The matter is plain
enough. A. Yes, the matter is that if we let formal identity hold for the
predicate, we deny it in each subject; whereas if we let it hold for each
subject, we deny it in the predicate. C. How so? A. If we posit
unconditionally ``animal is animal,'' then the horse becomes a fly and the fly
a snail; if we posit unconditionally ``the horse is horse, the fly is fly, the
snail is snail,'' then animal (horse) is not animal (fly) --- that is, the
common predicate contradicts itself and must fall away. C. This difficulty can
be overcome well enough. A. With dialectic or without? C. Without any dialectic
at all! One need only bethink oneself of the factual characteristic marks of
the concepts of existence, and of the relation of the disjunctive judgment to
\textit{diversi respectus}. A. That we must certainly hear! C. If we compare
animals with plants, it is seen that all existing animal individuals have one
thing in common: being animal. If on the other hand we compare the animal
existences with one another, then one animal series shows itself, by given
characteristic marks, different from another. If we compare the existences in
the one series with the existences in the other, then the individuals within
each series are alike; but
% --- p. 29 ---
if instead we compare the existences belonging to a given series with one
another, particular marks will show how one class is distinguished from
another, and so forth. The animal, then, in opposition to the plant, is what
the animal is; but by general and constant marks of difference the animal, in
opposition to itself, is either mollusc or articulate or vertebrate.\tnote{%
\emph{Bløddyr}, \emph{Leddyr}, \emph{Hvirveldyr}, and \emph{Straaledyr} just
below, are the Danish names of Cuvier's four \textit{embranchements}:
Mollusca, Articulata, Vertebrata and Radiata. The classification was standard
in 1861 and has since been abandoned; ``radiate'' covered what we should now
distribute among the echinoderms, cnidarians and others.} In opposition to
articulates and vertebrates the mollusc is always mollusc; but in opposition to
itself the mollusc is either a mollusc proper or a radiate. Now if one simply
goes on in this way with the divisions (\textit{divisiones} and
\textit{subdivisiones}) through every series, every class, every order, every
family, guided by the marks established by experience, until by way of genera
and species one has reached individuals --- what contradiction is there in
that, and what have such classifications to do with dialectic? A. Unless it be
that dialectic finds a good deal to do with \textit{diversi respectus}.
C. There is surely no danger of that. A. Then tell us, on the strength of your
completed classification: what actually is a horse? C. A horse (\textit{Equus
Caballus}) is, by the category of its kingdom, an animal; by the category of
its series, a vertebrate; by the category of its class, a mammal; by the
category of its order, a pachyderm; by its family category, akin to asses and
mules; and a species-category of its own within a genus-category of its own.
A. A fine collection of categories, then, all of which had to combine here to
yield the horse's concrete category of existence! If only one could now, with
the help of \textit{diversi respectus}, get it properly explained how these
many different categories have actually grown together --- not to say, been
bundled together --- into a single one! C. They of course enclose one another
in such a way that the more general always encompasses the more special.
A. Surely not, though, in the way one envelope encloses another, or one ring
encircles another? C. No, they rest upon one another, the more special having
in each case settled upon
% (p. 29 ends mid-sentence; the sentence continues on p. 30)

% II. Philosophy and the Particular Sciences begins p. 30 (General Remarks;
% A. Philosophy's Relation to Mathematics from p. 39).
% --- p. 30 ---
% (p. 30 opens with the last lines of part I's dialogue; the chapter head
% for II follows mid-page, after a section rule)
the more general. A. The more general a category is, the more abstract it is;
the more abstract it is, the further it is from existing: so the existent rests
upon the non-existent? C. No, I meant to say: they rest upon one another in
such a way that the more general is always borne by a more special, this by one
more special still, until we come to the most special of all, which is one with
the existent itself. A. So it is not being a horse that depends on being an
animal, but being an animal that depends on being a horse? C. No, that is not
how I conceived it either; indeed, at this moment I cannot, to speak plainly,
say how I actually did conceive it.

% (section rule as printed)

%% ============================================================
%%  II. PHILOSOPHY AND THE PARTICULAR SCIENCES  (pp. 30--476)
%% ============================================================
\chapter*{II.\\ Philosophy and the Particular Sciences.}
\addcontentsline{toc}{chapter}{II. Philosophy and the Particular Sciences}

% (short rule between chapter head and section head, as printed)
\section*{General Remarks.}
\addcontentsline{toc}{section}{General Remarks}

% (block quotation set letterspaced throughout, as printed; the citation
% is set unspaced)
\emph{``With respect to the \textit{a posteriori}, the advance of knowledge
will doubtless depend on new information's being constantly gathered, new
matters of fact tested, new discoveries made, and so on; but is the matter
thereby settled? If the advance of knowledge in relation to the \textit{a
posteriori} is regarded from this side alone, it is easy to dismiss philosophy.
If it cannot so much as exhaust the \textit{a priori} --- seeing that
mathematics owes it no thanks --- how much less can it meddle with positive
matters of fact! Let the historian lay a problem before it: it cannot answer,
for a historical problem is answered
% --- p. 31 ---
only historically. Let a chemist apply to it: it has no
% (small baseline speck between "intet" and "Greb" in the Danish; see the
%  transcription's note at this site)
knack for experiment and makes no discoveries in chemistry. Let the astronomer
want help with an observation: it scarcely understands him. Let the physician
ask it for advice: it is ignorant. Even what it has speculated out on its own
account concerning the relation of soul and body the physiologist may
unfortunately be unable to use.''} (Propædeut.\ pp.~99--103.)

The extent to which it holds that philosophy is dialectic is seen from its
relation to the subject-sciences,\tnote{\emph{Fagvidenskab}, literally the
science of a \emph{Fag} --- a subject, a trade, a department.
``Subject-science'' is used here, as in the translations of Nielsen's other
works, and its practitioner, the \emph{Fagmand}, is ``the specialist.'' These
are to be kept distinct from \emph{de særskilte Videnskaber}, ``the particular
sciences,'' of this part's title.} for that relation is dialectical in the
strictest sense. If one overlooks the dialectical, the relation becomes
ambiguous and doubtful; for while the subject-sciences are by nature inclined
to dismiss philosophy, philosophy in its turn, when it conceives its vocation
dogmatically, is tempted from its side to exalt itself at the expense of the
other sciences. Schelling's words are characteristic in this regard:
``Speculation is everything; it is to behold in God, to contemplate that which
is. Science itself has worth only in so far as it is speculative, that is to
say, contemplation of God as he is.'' (Propædeut.\ §~17, pp.~81--83).

What has chiefly contributed in recent times to confusing the issue is the
circumstance that Hegel --- whose penetrating grasp of contradiction, of the
negative, and of the immanent method of development (Propæd.\ pp.~150--158)
must place him so high in the ranks of the dialecticians --- has nevertheless,
in his account of philosophy's relation to the subject-sciences, overlooked the
dialectic of the relation. Instead of acknowledging the scientific independence
of the particular sciences, Hegel closes his propaedeutic survey of the
encyclopaedia of the sciences (Propæd.\ pp.~84--85) with the declaration that
the encyclopaedia is properly ``an exposition of the general content of
philosophy; for whatever in the sciences
% --- p. 32 ---
is founded upon reason depends upon philosophy, whereas whatever there is in
them of what rests on arbitrary and external determinations, or, as it is
called, is positive and statutory, lies outside it.'' But to let philosophy's
superiority rest in this way directly upon the subjection of the
subject-sciences is to take the matter undialectically; to see philosophy's
independence grounded in its dependence upon the independent subject-sciences
is, by contrast, to grasp the relation dialectically.

The independence of the subject-sciences over against philosophy is evident
from the course of their development; for the methods and discoveries that mark
the great advances of these sciences are so far from being attributable to
philosophy that ``the history of the inductive sciences'' (Whewell) shows, on
the contrary, to what degree independent research and a procedure answering to
the nature of the case have depended on the researcher's being able to free
himself from the one-sided influence of philosophical speculation (Propæd.\
pp.~167--68). What in the subject-sciences ``is founded upon reason'' by no
means depends upon philosophy in particular. In the first place it is not
reason that owes its origin to philosophy, but on the contrary philosophy that
owes its origin to reason. And next --- and in this connection this is
decisive --- it must be emphasized in the strongest terms, not only for the
sake of the subject-sciences but above all for philosophy's own sake (Propæd.\
§~23), that there is a knowledge which is scientific in the strictest sense
although it is not philosophical; indeed a thinking whose scientific energy
depends precisely on its not being philosophical. For that in the empirical
knowledge of the empirical there is a certainty of sensation, a many-sidedness
of representation, an added gift of intuition which no philosophical glimpse of
the Idea can replace, is surely as undeniable as it is undeniable that there is
an exact and critical thinking not to be exchanged for the dialectical
(Mathem.\ og Dialekt.\ p.~3). The rational
% --- p. 33 ---
distinctiveness of the particular sciences is seen, too, from the fact that
each one severally presupposes a scientific talent of its own answering to it.

But if one now sets up the proposition that \emph{the subject-sciences are
independent of philosophy, whereas philosophy --- particularly as regards the
appropriation of positive content --- is dependent on the subject-sciences},
then this presupposition, true as it is in itself, may easily occasion a false
inference. If philosophy, it is said, is so dependent on the exact and positive
sciences in everything that concerns the exact, the empirical and the positive,
then this dependence must plainly be a sign of its lack of independence. Every
time philosophy comes into contact with some determinate positive subject
matter, be its content mathematical, mechanical, chemical or historical, the
philosopher must of course put up with being not merely set right about the
facts but even duly schoolmastered by the specialist in the development of the
concepts themselves. It looks, then, as though the concepts were perfectly
trustworthy, clear, distinct and secure in all their consequences so long as
they are preserved in the form in which the subject uses them, but abandoned to
distortion as soon as they are made an object of philosophical reflection; it
looks as though the philosopher had to borrow everything from the specialist,
while the specialist for his part may regard the philosophical as a matter of
complete indifference. This way of seeing things rests on a misunderstanding.
Philosophy's dependence on the subject-sciences is, essentially regarded, a
sign of its independence, of its ideal self-validity.

In contrast to the material objects of the surrounding world all sciences are
of course ideal, all have ideality; but in contrast to philosophy, as the
representative of ideality, the subject-sciences must be said to stand for
reality. If the opposition is determined categorically, the real in its
finitude will be recognized by its asserting itself through excluding its
other, whereas the ideal proves its reality precisely by
% (signature "3" at the foot of p. 33)
% --- p. 34 ---
\setcounter{footnote}{0}
taking the other up into itself and relating itself to itself through the
other. That mathematics as mathematics, chemistry as chemistry, natural history
as natural history each keep to their
own\footnote{``There is a cardinal rule which can never be violated with
impunity, and which is nevertheless constantly violated in our groping advance
toward the light. It is: \emph{never to attempt the solution of the problems of
one science by the order of ideas belonging to another.} There is one class of
conceptions belonging to physics, another belonging to chemistry, a third
belonging to physiology, a fourth belonging to psychology, a fifth belonging to
social science. While all these sciences are closely connected, each has its
own independent province, which must be respected.'' G.~H. Lewes:
Hverdagslivets Physiologie, Danish translation, part 1, p.~66.\tnote{Lewes
wrote in English --- \emph{The Physiology of Common Life} (1859--60) --- and
Nielsen quotes the Danish version of it. The passage is therefore rendered back
into English here, and will not correspond word for word with Lewes's own
sentences.}}, and, even when contributions are fetched from other subjects,
still have first and directly to do each with its own, is precisely a mark of
these sciences' finite reality. That philosophy, by contrast, must always have
its ideal wealth in real poverty, its positive in its negative, was illustrated
by Plato in an ingenious myth: at Agathon's banquet Plato has his Socrates, as
is well known, present Eros, the ideal philosopher, as the son of Poros and
Penia.

The realistic science, once founded, may in the course of time need manifold
reshaping and undergo many changes, yet its supports and foundations will
through it all remain unshaken. If the idealistic science, on the other hand,
is to remodel its building, a \textit{de omnibus dubitandum} will gradually
shake every support; the whole building must then be pulled down to the ground,
that it may be raised in its new style from the ground up. Each time such a
critical juncture arrives it does indeed look for the moment as though reality
were at stake, as though it were now all over with philosophy; but history
confirms what also lies in the nature of the case, that such pervasive crises
mark precisely the turning points in philosophy's rebirth.

% --- p. 35 ---
% (block quotation letterspaced as printed; citation unspaced)
\emph{``The subject-sciences are solid; they always have steady work, ample
means and secure recognition in the republic of letters. Philosophy, by
comparison, is like a traveller without settled provision, a hovering shape
that sinks and rises, dies away and lives again, as the spirit will.''}
(Propæd.\ p.~102). --- Yet it is only from a wholly external, realistic
standpoint that the ideal is taken as merely the vague, the indeterminate, the
in itself indeterminable; the dialectical form is a form determinate and
determinable in itself, and philosophy's relation to the subject-sciences is
therefore to be stated categorically.

If philosophy is to be able to place itself in an inner and essential relation
to a particular science, it must discover in that science a problem which it
can also call its own problem. In order to show how such a \emph{common
problem} is chosen and treated, and so as not to scatter our attention over too
many points, we shall in what follows confine ourselves to examples drawn from
mathematics, chemistry and physiology; and we have thus, in mathematics,
\emph{the problem of infinity}; in chemistry, \emph{the problem of change}; in
physiology, \emph{the problem of life}.

But even if such common problems are characteristic enough when viewed from the
side of ideality, it is by no means thereby settled that they show themselves
so when they are grasped realistically. Because the philosopher has noted some
pervasive problem in one or another determinate subject-science, it is not yet
said that the specialist will acknowledge the same problem in a corresponding
way. A chemist who wished absolutely to resist a philosophical conception
might, for example, reason thus: ``There is in chemistry properly nothing at
all that one may call a problem of change, but rather a concept of change,
which arises and comes to be applied in so far as
% (signature "3*" at the foot of p. 35)
% --- p. 36 ---
chemistry investigates the natural laws according to which the inner changes of
bodies take place. But it is one thing to investigate the laws according to
which the changes take place, and another to want to fathom the essence of
change, to speculate about its possibility and impossibility --- a task wholly
foreign to chemistry.'' In the same spirit a physiologist wishing to dismiss
philosophy might reason thus: ``Physiology does indeed proceed from the
presupposition that there is a something we must call life, since in living
organisms there is a law-governed system of vital phenomena; but since
physiology essentially confines itself to investigating the phenomena, without
penetrating into what lies hidden behind them, it can by no means be granted
that there is here any proper problem of life. What life is in itself --- a
principle, a force, a product of the peculiar cooperation of mechanical and
chemical forces --- no matter! Let each form whatever individual representation
of the essence of life he pleases; so long as he is not thereby kept from
seeing the validity of the laws and demonstrating the physical causes,
physiology is satisfied.''

The difficulties which the natural influence of differing standpoints on the
conception of a scientific problem always brings with it are considerably
heightened by the ease with which something human --- that is to say, something
one-sided, accidental, arbitrary, passionate --- mixes itself into the
investigations. And since philosophy, for all its ideality, is like every other
science subject to human conditions, it becomes important in this connection to
take the matter from the purely human side as well, and so to note the
obligations and rights which the sciences have to uphold in their mutual
conflicts.

That the common problems in question must look otherwise in philosophy than in
the subject-science is a matter of course. The problem of change, even if the
chemist will acknowledge it in his own way, is in chemistry
% --- p. 37 ---
by no means what it is in philosophy; likewise the problem of life, even when
the physiologist grants its reality, is one thing in physiology and another in
psychology. The philosopher who knows the problem in the form of otherness, and
is aware that the same is nevertheless not the same, must therefore make it his
chief business to clarify the conditions of the development and to make the
transformation plain. Anyone who would seek guidance in a particular science
must of course make himself at home in its way of thinking. If the philosopher
neglects this --- if he misunderstands the character of the science in question
either as a whole or in essential particulars, whether by ascribing to it an
aim it does not have and cannot have, or by tearing its doctrines out of their
proper connection and so distorting them, or by overlooking the peculiar
presuppositions and conditions under which a given concept with its
corresponding designation was formed, thereby laying a false meaning under the
concept and in this way coming to force upon the science views, results and
consequences it cannot acknowledge --- then the specialist is within his rights
in speaking out. This right becomes a formal duty to the truth when such
blunders on philosophy's part, as has not seldom been the case (Propæd.\
pp.~170--171), are combined into the bargain with the conceit that it is
philosophy which from its exalted standpoint must feel called upon to reform
the subject-sciences. But further than a rejection of philosophical
encroachments the specialist cannot and should not go. If he goes further ---
so that, for example, he is not content to secure to the concept the meaning it
has in his science, but regards it as belonging to his subject to such a degree
that he either flatly forbids the philosopher to touch it, or at any rate takes
it upon himself to prescribe how far, in what manner and in what connections
philosophy may develop the borrowed concept graciously made over to it by
% --- p. 38 ---
the subject-science --- then there is to be seen in this not merely a
presumption but a ludicrous misunderstanding.

% (block quotation letterspaced as printed; citation unspaced)
\emph{``Confined to the several compartments of consciousness, the concepts
gained are still so grown together with finite constituents, still so weighed
down with particular presuppositions, that they find themselves as it were in a
state of unfreedom. They are held here as the prisoners of their matter within
a determinate horizon; they are, so to speak, the subject's bondsmen, who dare
not leave the appointed place of work.''} (Propædeutik p.~113.)

Now in setting out to free the real concepts from what is accidental in their
delimitation, to unfold the consequences and to bring the schematic forms into
dialectical flow, philosophy must
% (raised speck after "den" in the Danish; see the transcription's note)
of course --- once it has drawn the concepts into its own horizon, posed the
problems in its own way, and thus taken the whole responsibility of the
treatment upon itself --- have an indisputable right within its own province to
apply its own measure, to introduce its own designations and usage, and to
carry out the development in its own spirit. If it stands firm that every
subject-science must in essentials contain its own correctives, then philosophy
too must surely contain its correctives. A judgment upon the errors and
blunders a philosopher may possibly have committed in carrying through the
dialectical analyses belongs in the last instance not before the forum of the
subject-science but before that of philosophy.

To deny philosophy the full scientific right to treat in a free philosophical
manner what it has appropriated from the subject-sciences is to deny it a
development of concepts of its own, and to resort to a dictum of authority. If
it is to go by dicta of authority, then a specialist may settle philosophy's
fate as easily as that caliph settled the fate of the Alexandrian library:
``Either it contains only what stands in the Koran, and then it is superfluous,
% --- p. 39 ---
or it contains something else, and then it is false''; --- ``either philosophy
contains only what stands in the subject-science, and then it is superfluous,
or it contains something else, and then it is false.'' It is a Saracen dilemma;
but a Saracen dilemma cannot make philosophy give up its right.

% (section rule as printed)

\section*{A.\\ Philosophy's Relation to Mathematics:\\
  The Problem of Infinity.}
\addcontentsline{toc}{section}{A. Philosophy's Relation to Mathematics:
  The Problem of Infinity}

% (rule under the section head, as printed; the block quotation is
% letterspaced as printed, and here — unlike pp. 31, 35 and 38 — the
% citation is letterspaced too, so it stays inside the \emph)
\emph{``A philosophical fox or marten has stolen in upon mathematics' fat
poultry to fill his schematic belly; he means to suck out their warm blood, so
that nothing is left but the dead carcasses, void of sap and strength, which
are then to be presented as mathematics' vigorous, practical, lively concepts''
(Professor Steen's ``Afregning,'' p.~7).}\tnote{This whole section is a reply in
an ongoing quarrel. In 1860 the mathematician Adolph Steen attacked Nielsen's
\emph{Philosophie og Mathematik} in a pamphlet called \emph{Afregning}
(``Settling of Accounts''), and Nielsen answered with two pieces of his own ---
\emph{Hr.\ Professor Steens ``Afregning'': Hermed en specificeret Opgjørelse til
Revision} and \emph{Om Begrebet Nøiagtighed med særligt Hensyn paa Hr.\ Prof.\
Steens ``Bidrag''}, the latter a postscript to \emph{Philosophie og Mathematik}.
The short titles Nielsen cites here and below --- \emph{Afregning},
\emph{Begrebet: Nøiagtighed}, \emph{Om Begr.\ Nøiagt.} --- are to those pieces,
and the words he puts in quotation marks are his opponent's objections as he
reports them there. Without this the pages read as an argument with nobody.}

These boding --- indeed ill-boding --- words may well seem at first glance to
have something particularly forbidding about them, especially when set beside a
declaration such as this: \emph{``In the first place there is nothing that one
may name `the problem of infinity' in mathematics''} (cf.\ Begrebet:
Nøiagtighed, p.~38). But when it is thereupon conceded that there is \emph{``a
mathematical concept `infinite' which arises and comes to be applied, and which
is extremely important and moreover difficult, because it is of a mixed
nature''}, the forbidding element disappears and something inviting shows
itself. For if the mathematical concept ``infinite'' which arises and comes to
be applied is, on the mathematician's own admission, \emph{``extremely
important and moreover difficult, because it is of a mixed nature''}, then this
contains, so far as philosophy is concerned, a
% --- p. 40 ---
well-founded invitation to investigate whether this concept of infinity, in
mathematics so \emph{extremely important and moreover difficult}, may not after
% (the line-end of the corresponding Danish word may also be lightly
%  letterspaced; see the transcription's note at this site)
all be connected in one way or another --- perhaps without the mathematician's
knowing it --- with a problem running through the whole of mathematics; whether
philosophy may not have every reason to acknowledge, in this pervasive problem,
its own age-old problem of infinity; and whether mathematics, without a
mathematician's always quite thinking of it, may not be said to furnish
essential and indispensable contributions toward that problem's solution.

\subsection*{a) The Concept of Infinity.}
\addcontentsline{toc}{subsection}{a) The Concept of Infinity}

That the concept ``infinite'' which in mathematics ``arises and comes to be
applied'' is ``extremely important'' we may grant at once; but we must most
definitely deny that the concept of infinity, taken as a concept by itself, is
either difficult or of a mixed nature. Of the infinite in general only this can
be said: that it is a negation of the finite, the not-finite. If any difficulty
lay hidden in such a negation, the difficulty could not stem immediately from
the negation itself, but would have to be referred essentially to what the
negation cancels. So if one is to find difficulties in the mathematical
``infinite,'' one must find corresponding difficulties in the mathematical
``finite''; and if one raises no objection to the latter, it betrays mere
thoughtlessness to raise objections to the former.\tnote{Danish \emph{gjøre
Ophævelser over} is an idiom, ``to make a fuss about''; but \emph{Ophævelse} is
also the technical word for sublation, and Nielsen has just used its verb
(\emph{ophævet}, ``cancelled'') two sentences earlier. The pun does not
survive.} If the concept of infinity is difficult ``because it is of a mixed
nature,'' would not the concept of finitude, when the finite is specially
emphasized, be of a still more mixed nature? If it be objected that the
mathematical ``finite,'' which holds only of magnitudes and not of things,
cannot possibly be of a mixed nature, then one will also be obliged
% --- p. 41 ---
to grant that the negations in which such a ``finite'' disappears can as little
be so.

Concepts of a mixed nature always have a questionable admixture of the
conceptless, and as mixture-concepts they stand on the lowest rung in the order
of concepts. How does a mixture-concept come about? By parts of unlike kind,
brought together in an external and accidental way, being subsumed under a
collective unity, which in consequence comes to contain marks of unlike kind
accidentally brought together. Thus the concepts ``topsoil,'' ``gravel,''
``marl'' are genuine mixture-concepts; but a concept of gravel and a
mathematical concept of infinity --- what a difference!

But, someone may perhaps ask, if the concept ``infinite,'' in however many and
various connections it occurs and however many transformations it undergoes,
always preserves one and the same unalterable fundamental meaning --- namely
that of being ``the negation of finitude'' --- how does it come about that so
pure, so negative a logical concept can in a mathematician's eyes take on the
appearance of being ``of a mixed nature,'' or, what is the same, of belonging
among the mixture-concepts? The answer is not difficult: \emph{one mixes the
accidental with the essential; by this mixing one turns up a number of meanings
of differing kinds; one does not derive these separate meanings from the
fundamental meaning, does not exhibit the necessary transition of the one into
the other, but sets them up externally side by side, makes them an object of
enumeration --- and behold, the thing goes as of itself.} If one puts such
meanings as ``negation in positive form,'' ``impossible,'' ``transition
involving something qualitative,'' and so on, externally together, one gets at
once an impression of the mixed. How is one to distinguish a negation in
positive form from a negation in negative form, except by the former's
expressing the negative mediately and
% --- p. 42 ---
the latter immediately? But if now the finite is set against the infinite, it
is not the infinite that expresses the negation in positive form but on the
contrary the finite, in so far as the finite, though it has a positive form, is
yet mediately determined as the not-infinite. To take negation's own
immediately negative expression for its positive form betrays a comical
confusion. If someone, to display his acumen in matters logical, were to
maintain that the proposition ``$a$ is not something'' was straightforwardly
negative, whereas the proposition ``$a$ is nothing'' was affirmative, the
negation being lodged in the word ``nothing,'' one would surely agree without
difficulty to assign such trials of acumen --- formulated in all clumsiness
after a formal logic's \textit{S} is \textit{Non} --- \textit{P} --- if not
outright to the pedantic, then at any rate to that class of immature attempts
which betrays that the acute man has not yet learned to distinguish between
fruitful and unfruitful distinctions. If one now grants, so as not to be
saddled with too much logical mixed goods, that the proposition ``$a$ is
nothing'' is just as directly negative as ``$a$ is not something,'' then one
will also grant that the proposition ``$a$ is zero'' is just as directly
negative as ``$a$ is nothing.'' But if it stands firm that the sign $0$ is just
as direct and immediate an expression of negation as the word ``nothing,'' then
it takes a quite abnormal acumen to see that, while the proposition ``between
$a$ and $a$ there is no difference'' is straightforwardly negative, the
equation $a - a = 0$ should on the contrary be so nailed fast in the positive
that only the screw of a paraphrase could bore through to the negation. (Cf.\
Om Begrebet: Nøiagtighed, p.~39).

The logical construing of mathematical equations is a peculiar business; the
more weight is laid here on the conceptual side, the more evident it becomes
that philosophy has a word to say here too. On the natural mathematical way of
thinking, the equation $\frac{a}{0} = \infty$ shows us the well-known relation
between
% --- p. 43 ---
divisor and quotient in its extreme consequence: if, on the assumption of a
constant dividend, the divisor goes on shrinking to infinity, the quotient must
go on growing to infinity. There is nothing in this relation to recall the
impossible; everything points on the contrary to the possible: to the smallest
possible divisor there answers the greatest possible quotient. If, on the other
hand, in order that the infinite may come to signify the impossible, one
chooses another point of view and construes the equation thus --- ``When it is
asked how many times $0$ must be made greater in order to give $a$, the answer
in words is: that cannot be produced by multiplying $0$; but in signs:
$\frac{a}{0} = \infty$, $0$ must be made infinitely many times greater in order
to give $a$'' (Om Begr.\ Nøiagt.\ p.~40) --- then such a construal is well
fitted to promote all manner of admixtures of an illogical miscellany. If the
equation $\frac{a}{0} = \infty$ were really and seriously designed to answer
the in itself absurd question ``how many times must $0$, the absolute $0$, pure
nothing, be made greater in order to give $a$?'', then notations such as
$\frac{a}{0} = 0$, $\frac{a}{0} = a$ would have been far more adequate. For
from such expressions the impossible, the absurdity in the demand, is seen
instantly, whereas the equation $\frac{a}{0} = \infty$, by its ironic pointing
to the infinite, only mimically hints at what is absurd in the question, and
then --- quite rightly --- transfixes the questioner by ``a negation in
positive form,'' because he is asking not as a mathematical thinker but as a
logical fool. If one gives up the logical foolery of trying to reckon how many
times pure zero must be made greater in order to produce $a$; if, in order that
the equation may yield any sense, one lets $0$ denote the smallest possible
divisor while $\infty$ denotes the corresponding greatest possible quotient ---
then it is seen from every side
% --- p. 44 ---
how the impossible gives way to the possible. From the equation $\frac{a}{0} =
\infty$ there follows directly, in virtue of the concept of the relation in
question, $a = \infty \frac{a}{\infty} = \infty.0$; likewise from
$\frac{b}{0} = \infty$, $\frac{c}{0} = \infty$, $\frac{d}{0} = \infty\ldots$
there follow $b = \infty.\frac{b}{\infty} = \infty.0$,
$c = \infty \frac{c}{\infty} = \infty.0$,
$d = \infty.\frac{d}{\infty} = \infty.0\ldots$
% (the multiplication dot after the first ∞ is printed inconsistently in the
% Danish: present in the b- and d-equations and at "altsaa: a = ∞.a/∞" below,
% absent in the a- and c-equations; reproduced as printed)
Instead of the impossibility of $\infty.0$'s being able to give $a$, we see on
the contrary the possibility of $\infty.0$'s giving not only $a$ but $b$, $c$,
$d\ldots$ --- in short, as many different magnitudes as you please. Instead of
the curt, dismissive ``impossible,'' we get possibility upon possibility,
possibilities in abundance. This result is of course so far from surprising the
\emph{thinking} mathematician that it does not even surprise ordinary sound
human understanding. Everyone understands that any finite magnitude whatever
may --- and this is a possibility --- be said to consist of infinitely many
infinitely small parts, hence $a = \infty.\frac{a}{\infty} = \infty.0$, and
that it may therefore also --- and this is a possibility --- be divided into
infinitely many infinitely small parts, hence $\frac{a}{\infty} = 0$.

The fixed idea that the concept ``infinite,'' as a ``negation in positive
form,'' should be better fitted to denote ``the impossible'' than the concept
``finite'' is easy to clear up. For ``the impossible'' is in and for itself as
indifferent to ``infinite'' as to ``finite''; the impossible shows itself every
time a demand that has been set up annihilates itself, and such
self-annihilating demands may just as well presuppose the finite as the
infinite. It is, to stay with the infinite, impossible to \emph{enumerate} all
the terms of an infinite series; it is, to pass to the finite,
% (p. 44 ends mid-sentence in the Danish; the sentence continues on p. 45)

% --- p. 45 ---
impossible for a human being to \emph{count up to} a quadrillion; and it is ---
here we must lean at once on both the finite and the infinite --- impossible
that a sphere with an infinite radius should be a sphere, or a plane a
spherical surface with a finite radius; and just as impossible that the
equation $a = 2a = \infty.a = \frac{0}{a}$ should have any reality. If in the
language of mathematical signs one wants a more direct notation for something
impossible in itself, then it is not $0$ and $\infty$ but $\sqrt{-1}$,
$\sqrt[2m]{-a}$ and the like that one has to keep to. (The ``Afregning'' as
revised by R.~N. Appendix, pp.~26--27).

In logical generality it is easily seen that the concept ``infinite'' does not
as a concept take up constituents of unlike kind, and so cannot be reckoned
among the mixture-concepts. That the concept nevertheless occurs in several
meanings, shows itself in a different light, and to that extent may, when
treated superficially, acquire a tinge of being ``of a mixed nature,'' has its
ground in something quite different: to think the concept through is to unfold
the problem.

% Run-in head, not a sectioning command in the Danish; the \addcontentsline is
% supplied so the Indhold's entry has a counterpart in the contents.
\textbf{b) The Problem of Infinity.}
\addcontentsline{toc}{subsection}{b) The Problem of Infinity}
Suppose someone were to open an objection to a philosophical reference to the
way in which the problem of finitude finds its solution in mathematics with the
following declaration: ``In the first place there is nothing that one may name
the problem of finitude in mathematics, but rather a mathematical concept
`finite' which arises and comes to be applied, and which is extremely important
and moreover difficult, because it is of
% the Danish print mis-sets the marks here: the objection opens „Der ... and the
% compositor closes the quote both after "Endelighedsproblemet" and after
% "Natur" without re-opening one — one “ in excess on p. 45, offset by the
% unclosed „ on p. 46. Neither defect can be reproduced in English; both places
% are set normally here and logged, as in the transcription.
a mixed nature.'' One would surely find that a strange thing to say. What a
finite magnitude is, is soon said; but the problems concerning the mutual
relations of finite magnitudes are countless. Instead of saying that there is
no problem of finitude in mathematics, one would rather have to say
% --- p. 46 ---
that there are countless problems of finitude in mathematics. If these problems
fell apart from one another in loose multiplicity, one might indeed deny that
there was any general problem running through them all; but one would then at
the same stroke have to deny that there was any mathematical method --- in
other words, that mathematics was a science. If mathematics is a science, if
there is a method by which the various problems of finitude are framed and
solved in systematic connection, then there is also a general problem running
through them all, a fundamental problem embracing and determining the method in
all its ramifications, a problem whose all-round solution demands nothing less
than the whole of ``the analysis of the finite.''

But what holds so incontestably of ``the problem of finitude'' in its relation
to ``the analysis of the finite'' --- should this not hold also of the relation
between ``the problem of infinity'' and ``the analysis of the infinite''? What
holds so incontestably of the concept ``finite'' in its relation to ``the
problem of finitude'' --- namely that the difficulty lies not in the concept
but in the problem --- should this not, on closer inspection, be found to hold
also of the concept
% the quote opened here before „Uendelighedsproblemet?“ is never closed in the
% Danish print (balance +1); set normally in English.
``infinite'' in its relation to ``the problem of infinity''? ---

\emph{α) Concept and Problem.} As the distinguishing mark of mathematical
concepts --- however much they may otherwise resemble the general logical ones
--- it must be expressly emphasized (Mathem.\ og Dial.\ pp.~21--27) that they
are operational concepts. From this it follows that, while in their generality
they present no difficulty, analysing thought will soon discover that their
proper development must always take place in the closest connection with the
operations whose logical essence they are meant to clarify. What a logarithm is
can easily be stated in a definition exact enough for sound understanding to
hold the concept fast; but the manifold and characteristic consequences in
which the concept's content first really
% --- p. 47 ---
shows itself come forth only in the course of the operations. What is meant by
dividing a line in extreme and mean ratio --- or, as it is also called, by
golden-sectioning the line --- is soon said and easily grasped; the concept
``golden-section''\tnote{The Danish verb is \emph{høidele}, literally ``to
high-divide,'' the nineteenth-century Danish term for division in extreme and
mean ratio, i.e.\ the golden section. English has no corresponding verb, so
``golden-section'' does duty as one here and in the next sentence.} presents no
difficulty. The difficulty, if there is one, is to solve the problem of the
golden section and to develop all the consequences contained in it. What holds
of the particular operational concepts holds of the general ones too: only from
the systematic solution of the problems of finitude does it come to light what
mathematically lies in the concept ``finite''; only a carrying-through of ``the
analysis of the infinite'' can open us a scientific insight into the
mathematical concept of infinity. From the consequences of the operations it is
further seen that the infinite falls not only outside but also inside the
finite (cf.\ Phil.\ og Mathem.\ pp.~20--24; Mathem.\ og Dialekt.\ pp.~28--38),
and that it develops in manifold ways out of finitude itself. Had a pervasive
problem of infinity not originally stirred beneath ``the analysis of the
finite'' and declared itself in many separate flashes, the differential and
integral calculus would never have been developed. When Leibniz solved ``the
problem of tangents,'' what was essentially at stake was not the tangent and
the secant (cf.\ Phil.\ og Mathem.\ p.~64) but the articulating reciprocity of
finite and infinite: in ``the problem of tangents'' ``the problem of infinity''
is to be seen.

\emph{β) The Mathematical Conception.} That thinking mathematicians, however
much weight they may otherwise lay on the satisfactory results of the
operational method, have nevertheless never been able to separate
satisfactorily the concept they find inconceivable from the method that rests
upon that concept, is shown by the history of mathematics from Leibniz and
Newton (cf.\ Mathem.\ og Dial.\ pp.~90--91, p.~95 and elsewhere) down to the
present time. ``In the first place, when the principle is unknown we cannot
know whether all the consequences following from it can be admitted: we go
% --- p. 48 ---
\setcounter{footnote}{0}
as if blindfold, and it might well happen that after laborious effort we should
at last strike upon something which exhibited precisely the nature of the
principle --- that is to say, upon something we could not conceive. In the
second place: is it not a shameful degradation of reason to make it follow a
concatenation of terms, a series whose beginning and end it cannot see through?
It is a great evil that is born of that kind of routine: accustomed to
operating without seeing the ground, the mind loses all practice in thinking,
and ends by carrying everything out quite mechanically.'' (Mathem.\ og Dial.\
p.~III). These misgivings of Savrien's
% "Savriens" so printed in the Danish (presumably for Savérien); not corrected.
sufficiently explain the unrest and tension in which the analysis of the
infinite has constantly held mathematical
reflection\footnote{\textit{See Carnot: Reflexions sur la Métaphysique du
Calcul Infinitésimal. Paris 1797} (übersetzt von Hauf. Frankf. a.~M. 1800).
Versuch einer Philosophie der Mathematik --- mit einer Kritik der Aufstellungen
Hegel's über den Zweck und die Natur der höheren Analysis --- von H. Schwartz.
Halle 1853.} (Mathem.\ og Dial.\ pp.~55--56), whose problem --- as a prize
question set by the Berlin Academy in the year 1784 characteristically puts it,
``wie aus einer widersprechenden Voraussetzung so viele wahren Sätze haben
können hergeleitet werden'' [how so many true propositions could have been
derived from a contradictory presupposition] --- will not cease to return, even
though ``routine,'' the mechanical facility, will at all times no doubt afford
``\textit{habiles Calculateurs}'' satisfaction so ample that it may easily
escape this or that reckoning-master's notice that there is in mathematics
something one may name ``the problem of infinity.''

\emph{γ) The Philosophical Conception.} ``To make the problem of infinity
recognizable, reflection must lean upon the operation; to solve the problem, it
must appropriate the intelligence expressed in the operation. That there can be
a finite ratio between infinite magnitudes; that a sum of infinitely many terms
may in some cases make up a finite and
% --- p. 49 ---
in others an infinite magnitude; that a series whose terms are so related that
the limit of the ratio of two successive terms gives the quotient 1 may in some
cases be convergent and in others divergent, and so on --- such things will
everywhere escape the thinking that does not, at the operation's hint, catch
sight of the exact determinations. That these determinations must in
philosophical form become dialectical, while in mathematical form they neither
are nor can be so, shows a characteristic difference between mathematics and
philosophy, but by no means any conflict. No concept, however exactly it be
defined, can maintain itself in a being indifferent to the other concepts: the
validity of concepts is seen in their transitions; but the transitions which
exact analysis makes by the operation, the thinking that penetrates the
operational concepts makes by dialectic'' (Mathem.\ og Dial.\ p.~57. Cf.\
pp.~59--60). If the finite is set immediately against the infinite, sound
understanding involuntarily gets the impression that it is the finite which, in
its multiplicity of terms and relational determinations, contains the wealth of
existence, whereas the infinite is a bottomless nothing, waste and empty as
empty space. And since thought can determine only what is in itself
determinable, and only that is determinable which admits in some way of being
defined, and only that admits of being defined which admits of being bounded,
and the bounded is a finite --- it must constantly appear to the understanding
as though only the finite were conceivable, while the infinite, the unbounded,
were really not to be conceived at all, but, like the ancients' hyle
(τὸ ἄπειρον), to be known by a kind of spurious thinking
(λογισμὸς νόθος).\tnote{The phrase is Plato's, \emph{Timaeus} 52b, where the
receptacle of becoming is said to be apprehended by a ``bastard reasoning'' ---
$\lambda o\gamma\iota\sigma\mu\tilde{\omega}$ $\tau\iota\nu\iota$
$\nu\acute{o}\theta\omega$ --- and hardly to be believed in. Nielsen pairs it
with \emph{hyle}, matter, and with the Aristotelian
$\tau\grave{o}$ $\ddot{\alpha}\pi\epsilon\iota\rho o\nu$, the unlimited.} Where
this misunderstanding has taken hold there can of course be no question of
entering more closely into the problem of infinity; to see the problem's
philosophical significance one must absolutely have at least an inkling of the
dialectical reciprocity of the finite and the
% --- p. 50 ---
% degrees are printed as a superscript 0 (not a ring) here and on p. 59;
% reproduced as ^{0} throughout, as in the transcription.
infinite. ``An angle of $90^{0}$ can be divided into infinitely many infinitely
small angles, and so can an angle of $1^{0}$; a square mile can be divided into
infinitely many infinitely small squares, and so can a square inch: what
follows from this? If the emphasis is laid on the resolution of the finite into
the infinite, the consequence does indeed seem to be that there is no longer
any difference between the angle of $90^{0}$ and the angle of $1^{0}$, that the
part becomes equal to the whole --- that the boundary, in other words, is
effaced. But if the emphasis is laid instead on the infinite's existing in and
through the finite, then there is not only a difference between finite and
finite, between the part and the whole, but at the same stroke a definite
difference shows itself between infinite and infinite, with respect to the
infinitely small as well as to the infinitely great. The one infinite number of
infinitely small angles may then be as different from the other as $1^{0}$ from
$90^{0}$; the one infinite number of infinitely small squares as different from
the other as a square inch from a square mile. That the boundary is sublated in
the finite means that it is posited in the infinite; that the infinite takes
the boundary up into itself means that the finite is posited. The dialectical
unity shows itself in this, that the boundary belongs as much to the infinite
as to the finite'' (Mathm.\ og Dial.\ p.~33). The fundamental thought that the
infinite not only effaces the boundaries of finitude but also takes the
sublated boundary-determinations up as articulating moments in its own
development of concepts is, from the standpoint of the operations, sufficiently
designated by saying that the infinite is analysed.

% NB the printed head reads "b)." where the Indhold reads "c)" — see the
% transcription's INDHOLD COLLATION note.
\textbf{b). The Analysis of the Infinite.}
\addcontentsline{toc}{subsection}{b). The Analysis of the Infinite}
% [text to be added: pp. 50--80]

\section*{B.\\ Philosophy's Relation to Chemistry:\\
  The Problem of Change.}
\addcontentsline{toc}{section}{B. Philosophy's Relation to Chemistry:
  The Problem of Change}

% [text to be added: p. 81 -- the section's own intro, before a)]

\subsection*{a) Chemistry as a Distinct Science.}
\addcontentsline{toc}{subsection}{a) Chemistry as a Distinct Science}

% [text to be added: p. 82]

% The head of α) carries a footnote in the Danish (Leçons ...); carry it over.
\subsection*{α) The Origin of Chemistry.}
\addcontentsline{toc}{subsection}{α) The Origin of Chemistry}

% [text to be added: pp. 83--91]

\subsection*{β) Chemical Forces.}
\addcontentsline{toc}{subsection}{β) Chemical Forces}

% [text to be added: pp. 92--102]

% Printed "γ." with a period, against the Indhold's "γ)" — as transcribed.
\subsection*{γ. Chemical Combinations.}
\addcontentsline{toc}{subsection}{γ. Chemical Combinations}

% [text to be added: pp. 103--120]
% NB pp. 95 and 96 carry the double-decomposition schemes; use \decompscheme.

\subsection*{b) The Problem of Change in Chemistry.}
\addcontentsline{toc}{subsection}{b) The Problem of Change in Chemistry}

% [text to be added: pp. 121--123]

\subsection*{α) Atom and Equivalent.}
\addcontentsline{toc}{subsection}{α) Atom and Equivalent}

% [text to be added: pp. 124--129]

\subsection*{β) Atomic Constitution and Crystal Form.}
\addcontentsline{toc}{subsection}{β) Atomic Constitution and Crystal Form}

% [text to be added: pp. 130--134]

\subsection*{γ) Atom and Substance.}
\addcontentsline{toc}{subsection}{γ) Atom and Substance}

% [text to be added: pp. 135--142]

\subsection*{c) The Problem of Change in Philosophy.}
\addcontentsline{toc}{subsection}{c) The Problem of Change in Philosophy}

% [text to be added: pp. 143--144]

\subsection*{α) The Atom as Point of Discretion.}
\addcontentsline{toc}{subsection}{α) The Atom as Point of Discretion}

% [text to be added: pp. 145--150]

\subsection*{β) The Atom as Point of Function.}
\addcontentsline{toc}{subsection}{β) The Atom as Point of Function}

% [text to be added: pp. 151--159]

\subsection*{γ) The Atom as the Substantial Moment.}
\addcontentsline{toc}{subsection}{γ) The Atom as the Substantial Moment}

% [text to be added: pp. 160--166]

\section*{Result.}
\addcontentsline{toc}{section}{Result}

% [text to be added: pp. 167--176]

\section*{C.\\ Philosophy's Relation to Physiology:\\
  The Problem of Life.}
\addcontentsline{toc}{section}{C. Philosophy's Relation to Physiology:
  The Problem of Life}

% [text to be added: pp. 177--199]
% NB the opening paragraph of section C (p. 177) is set WHOLLY gesperrt in the
% print and is \emph{} throughout in the transcription; carry the emphasis over.
% This span is footnote-heavy (pp. 178, 179, 181, 184--187, 189, 191--194, 199),
% with several notes running over one or two following pages unmarked.

\subsection*{a) The Essence of Life.}
\addcontentsline{toc}{subsection}{a) The Essence of Life}

% [text to be added: pp. 200--207]

\subsection*{α) The Organism.}
\addcontentsline{toc}{subsection}{α) The Organism}

% [text to be added: pp. 208--235]

% Printed "β." with a period, against the Indhold's "β)" — as transcribed.
% The head stands at the FOOT of p. 236, so the text runs on from p. 236.
\subsection*{β. The Organic Type.}
\addcontentsline{toc}{subsection}{β. The Organic Type}

% [text to be added: pp. 236--260]

\subsection*{γ) The Organic Purpose.}
\addcontentsline{toc}{subsection}{γ) The Organic Purpose}

% [text to be added: pp. 261--285]

\subsection*{b) The Stages of Life.}
\addcontentsline{toc}{subsection}{b) The Stages of Life}

% [text to be added: pp. 286--290]

\subsection*{α) Plant Life.}
\addcontentsline{toc}{subsection}{α) Plant Life}

% [text to be added: pp. 291--327]
% NB pp. 326--328 are the Darwin passage, already translated as the excerpt
% edition texts/nielsen/propaedeutik-darwin/ — collate against it, and keep the
% terminology settled there (Paradigme = paradigm, Anlæg = disposition,
% Selvhed = selfhood, Andethed = otherness, Concretionsloven = law of concretion).

\subsection*{β) Animal Life.}
\addcontentsline{toc}{subsection}{β) Animal Life}

% [text to be added: pp. 328--398]

% Printed "γ." with a period, against the Indhold's "γ)" — as transcribed.
\subsection*{γ. Human Life.}
\addcontentsline{toc}{subsection}{γ. Human Life.}

% [text to be added: pp. 399--422]

\subsection*{c) Soul and Body.}
\addcontentsline{toc}{subsection}{c) Soul and Body.}

% [text to be added: p. 423]

\subsection*{α) Metaphysical Theories.}
\addcontentsline{toc}{subsection}{α) Metaphysical Theories.}

% [text to be added: pp. 424--436]
% NB p. 425 carries Descartes' Amsterdam imprint in apostrophus numerals,
% CIƆ IƆCLXXVII; the reversed C is declared in the preamble.
% NB p. 430 head "ββ) Identitetslæren." against the Indhold's "ββ) Identitetslære".

\subsection*{β) Psychophysical Investigations.}
\addcontentsline{toc}{subsection}{β) Psychophysical Investigations.}

% [text to be added: pp. 437--467]
% This is the section Nielsen says (Preface) was NOT delivered orally at all.

\subsection*{γ) The Limits of Human Knowledge.}
\addcontentsline{toc}{subsection}{γ) The Limits of Human Knowledge.}

% [text to be added: pp. 468--476]

%% ============================================================
%%  III. PHILOSOPHY AND ACADEMIC STUDY  (pp. 477--482)
%% ============================================================
\chapter*{III.\\ Philosophy and Academic Study.}
\addcontentsline{toc}{chapter}{III. Philosophy and Academic Study}

% The three subdivisions of Part III are set RUN-IN, in BOLD Fraktur (not
% letterspaced), at the head of their paragraph — so no sectioning command,
% exactly as with the run-in heads elsewhere. The \addcontentsline lines are
% supplied only so that the Indhold's three entries have counterparts in the
% table of contents.
\addcontentsline{toc}{section}{Scientific Thinking}
\addcontentsline{toc}{section}{The Connection of the Sciences}
\addcontentsline{toc}{section}{Scientific Cultivation}

% [text to be added: pp. 477--482]

%% ============================================================
%%  CONTENTS  (the book's own Indhold, printed at the end, p. 482 area)
%% ============================================================
\chapter*{Contents.}
\addcontentsline{toc}{chapter}{Contents (the book's own Indhold)}

% [text to be added: the book's own Indhold, 80 entries, PDF 11--12]
% Translate the entries to match the heads used above; keep the printed page
% numbers, which the transcription's INDHOLD COLLATION verified as all correct.
% Four Indhold entries have NO printed head in the body — "Muligheden og
% Modsigelsen" (8), "Modsigelsens Grundsætning" (11), "Anvendelse af
% Modsigelsens Grundsætning" (13), "Tvetydigt og Tvesidigt" (15). Record them,
% do not reconcile them.

\end{document}
